% Route B -- lifting transcendence from the relaxed series V to the true series U.
% All numbered facts validated in code/zeta_probe/route_b (uv_perelem.py, joint_W.json,
% defect_probe.py): the bridge U=W(x,1), V=W(x,x^{-2}) reproduces u_n EXACTLY for
% n=0..22 and v_n for m=0..14 (the latter limited only by BFS truncation).

\section{From the relaxed series to the true series: the cycle bridge}
\label{sec:lifting}

Throughout $q=x^{2}$.  Recall $u_n=|\{g:\ell_T(g)=n\}|$ (true word length) and
$v_n=|\{g:\ell_R(g)=n\}|$ (relaxed length), $U=\sum u_nx^n$, $V=\sum v_nx^n$.
The Metric Theorem gives, for every group element $g$,
\begin{equation}\label{eq:plus2}
  \ell_T(g)=\ell_R(g)+2\,c(g),\qquad
  c(g):=\#\{\text{isolated cycles in any relaxed-optimal realization of }g\}\in\mathbb Z_{\ge0}.
\end{equation}
(The $+2$ per isolated cycle is the cost of splicing a closed disconnected excursion
into the single Eulerian trail.)  Equation~\eqref{eq:plus2} is validated per element
against breadth-first true distance and the closed-form relaxed length: on all
elements of true length $\le 16$ the difference $\ell_T-\ell_R$ is even and
non-negative, with $c$ the half-difference (\texttt{uv\_perelem.py}: $0$ odd, $0$
negative).

\subsection{The two-variable bridge}

\begin{definition}
$\displaystyle W(x,y)=\sum_{n,c\ge0}w_{n,c}\,x^{n}y^{c}$, where
$w_{n,c}=|\{g:\ell_T(g)=n,\ c(g)=c\}|$.
\end{definition}

\begin{proposition}[Bridge]\label{prop:bridge}
$U(x)=W(x,1)$ and $V(x)=W(x,x^{-2})$.
\end{proposition}

\begin{proof}
$U(x)=\sum_n x^n\sum_c w_{n,c}=W(x,1)$ since $u_n=\sum_c w_{n,c}$.  By
\eqref{eq:plus2}, an element counted by $w_{n,c}$ has relaxed length $n-2c$, hence
$v_m=\sum_{n-2c=m}w_{n,c}$ and $V(x)=\sum_{n,c}w_{n,c}x^{n-2c}=W(x,x^{-2})$.
\end{proof}

\noindent\emph{Convention \textup(the symbol $W$\textup).} Here $W(x,y)=\sum_{n,c}w_{n,c}\,x^{n}y^{c}
=\sum_g x^{\ell_T(g)}y^{c(g)}$ is keyed on the \emph{true} length $\ell_T$, so $U=W(x,1)$ and
$V=W(x,x^{-2})$. In the transcendence section the same letter denotes the \emph{relaxed}-keyed
$\widehat W(x,y)=\sum_g x^{\rho(g)}y^{c(g)}=W(x,x^{-2}y)$ \textup(see~\eqref{eq:UWxx2}\textup), for which
$U=\widehat W(x,x^{2})$ and $V=\widehat W(x,1)$. The two are interchangeable via $\rho=\ell_T-2c$; no
statement depends on the choice.

\noindent\emph{Validation.} With $w_{n,c}$ tabulated from breadth-first enumeration
to radius $22$ (\texttt{joint\_W.json}), $W(x,1)$ reproduces $u_0,\dots,u_{22}$
\emph{exactly} (all $23$ terms), and $W(x,x^{-2})$ reproduces $v_0,\dots,v_{14}$
exactly (higher $m$ truncated only because elements of relaxed length $m$ but true
length $>22$ are absent from the radius-$22$ enumeration).  The defect
$D=V-U=\sum_{n,c}w_{n,c}\,(x^{n-2c}-x^n)$ involves only the $c\ge1$ part.

\subsection{The connectivity defect is bulk-localised}

Write the travel interval of $g=(\varepsilon^\ast,\delta^\ast,k^\ast;P)$ as
$I_k=[\min(0,k),\max(0,k))$, the edges carrying the travel indicator $f=\pm1$.

\begin{proposition}[Localisation of cycles]\label{prop:local}
\textup{(i)} If the lamp support of $g$ lies inside $I_k$ \textup(``pure travel''\textup),
then $c(g)=0$.  \textup{(ii)} Every $g$ with $c(g)\ge1$ has a lamp deposit on some
edge outside $I_k$.
\end{proposition}

\noindent\emph{Validation.} Over all elements of true length $\le16$: \textup{(i)}
$0$ violations out of $3869$ pure-travel elements; \textup{(ii)} all $234$
cycle-bearing elements carry an out-of-interval (bulk) deposit
(\texttt{uv\_perelem.py} structural scan).  Mechanistically, the travel interval is
the forced single $0\!\to\!k$ backbone of the Eulerian trail; an isolated cycle is a
closed bulk excursion ($f=0$ region) separated from that backbone by a reachability
gap, so it can only originate outside $I_k$.

\subsection{The travel block is connectivity-invariant}

\begin{proposition}[Invariance of the travel block]\label{prop:travelinv}
Let $T(x)=\sum_n t_n x^n$ count pure-travel elements by length.  Then the true and
relaxed gradings of $T$ coincide: the same series $T$ is obtained whether elements
are bucketed by $\ell_T$ or by $\ell_R$.  Consequently the travel section recursion
of \S\ref{sec:routeb} \textup(producing $A_k,C_k,\Sigma_k$\textup) is identical in
both models, and the cycle marker $y$ does not enter it: the travel block is
$y$-independent.
\end{proposition}

\begin{proof}
By Proposition~\ref{prop:local}(i) every pure-travel element has $c=0$, so
$\ell_T=\ell_R$ on this set and the two gradings of $T$ agree termwise.  The travel
section recursion of \S\ref{sec:routeb} is assembled solely from pure-travel runs
(appending one $f=\pm1$ edge at a time), so it is computed inside a $c\equiv0$
sector; in $W(x,y)$ the travel block carries $y^0$ only.
\end{proof}

\noindent\emph{Validation.} The pure-travel generating function is computed in both
gradings to radius $18$ and is bit-for-bit identical:
$1,3,5,8,12,17,24,35,52,78,116,173,256,383,572,858,1276,1907,2832$
(\texttt{uv\_perelem.py}); its growth ratio $\to\beta_2=1.49162$, the travel-pole
rate.

\subsection{Pole survival}

By Theorem~\ref{thm:main} (the relaxed result, modulo Lemma~\ref{lem:cos}), the
travel block $G^T_0=\Sigma_0/(1-\Sigma_1)$ has infinitely many poles
$q_1=q^\ast<q_2<\cdots\uparrow1$, the roots of $\Sigma_1(q)=1$.  Two further facts
hold unconditionally.

\begin{lemma}[Bulk block misses the travel poles]\label{lem:missing}
At each travel pole $q_m$ one has $1-S_1(q_m)\ne0$; i.e.\ the relaxed bulk block
$S_0/(1-S_1)$ is holomorphic at every $q_m$.  Moreover its nearest singularity is at
$q_b=0.609567\ldots>q^\ast$, so it is holomorphic on the disc $|q|<q_b$ containing
$q^\ast$.
\end{lemma}

\noindent\emph{Validation.} \texttt{final\_verify2.py} / scan: $1-S_1(q_m)=
0.520,-0.210,0.127,-0.058,0.049,-0.033,\dots$ at $q_1,\dots,q_6$ (never $0$); the
bulk and travel pole sets interlace.

\begin{theorem}[Conditional lifting]\label{thm:lift}
Assume the assembly factorisation \textup{(A)} below.  Then $U$ has infinitely many
poles in $|q|<1$ accumulating at $q=1$, hence $U$ is transcendental over $\mathbb
Q(x)$.
\end{theorem}

\noindent\emph{Update on the conditionality.} The pole \emph{accumulation} half --- that
$\{1-\Sigma_1=0\}$ is an infinite set accumulating at $q=1$ --- is the \emph{same} statement
needed for $V$, and is now established conditional only on the cited steepest-descent theorem
(transcendence.tex, Lemmas \texttt{lem:Bbounded}, \texttt{lem:extremephase}: the extreme-phase
bound $|E(m\pi)|<1$, with the amplitude's finite variation proved by the no-resonance Stirling
cancellation), \emph{not} on the open uniform Lemma~\ref{lem:cos}. By
Proposition~\ref{prop:travelinv} the travel poles $q_m$ are the identical, $y$-free object in
$U$ and $V$, so $U$ inherits this accumulation. Hence the \emph{only} remaining gap for $U$ is
(A): the assembled numerator $\mathcal N_U(q_m)\ne0$ at infinitely many $q_m$.

\paragraph{(A) Assembly factorisation (the remaining gap).}
The full generating function of either model is a sum over the displacement $k$,
$\;\mathrm{GF}=\sum_{k}\Theta_k(\text{bulk dressings})$, in which the travel run of
length $|k|$ contributes a transfer whose $k$-sum telescopes to the travel resolvent
$(1-\Sigma_1)^{-1}$ times a bulk-dressing factor $B_\bullet(q)$ that is holomorphic
at every $q_m$.  Statement (A) is that, in the \emph{true} model, the corresponding
bulk-dressing factor $B_U(q)$ is likewise holomorphic and does \emph{not vanish} at
the $q_m$, so the factor $(1-\Sigma_1)^{-1}$ -- which by
Proposition~\ref{prop:travelinv} is the identical, $y$-free object in both models --
survives in $U$ with non-zero residue.

\begin{proof}[Proof of Theorem~\ref{thm:lift} given (A)]
By Proposition~\ref{prop:travelinv} the denominator $1-\Sigma_1$ is the same in $U$
and $V$.  By (A), $U=B_U\cdot(1-\Sigma_1)^{-1}+(\text{terms regular at the }q_m)$
with $B_U(q_m)\ne0$, so each $q_m$ is a genuine pole of $U$.  As $\{q_m\}$ is
infinite and accumulates at $q=1$, $U$ has infinitely many singularities; a $D$-finite
function has finitely many, so $U$ is transcendental
(Theorem~\ref{thm:main}'s argument verbatim).  The substitution $q=x^2$ preserves
transcendence over $\mathbb Q(x)$.
\end{proof}

\begin{remark}[Status]\label{rem:liftstatus}
Propositions~\ref{prop:bridge}--\ref{prop:travelinv} and Lemma~\ref{lem:missing} are
rigorous (and exhaustively validated).  They reduce the transcendence of $U$ to the
single structural statement (A): \emph{the cycle-corrected bulk dressing $B_U$ is
holomorphic and non-vanishing at the travel poles $q_m$.}  The relaxed analogue of
(A) is Lemma~\ref{lem:missing} (with $B_V=S_0/(1-S_1)$, manifestly holomorphic and
generically non-zero at the $q_m$).  For $U$ the bulk dressing is the
\emph{connectivity-corrected} block $W^{\mathrm{bulk}}(x,1)$ rather than
$S_0/(1-S_1)$; we have shown it is built from the same bulk transfer with the
cycle-creating transitions removed, and that the cycle marker does not touch the
travel recursion, but we have not proved $B_U$ holomorphic up to $q_b$ nor
$B_U(q_m)\ne0$ in closed form.  Hence, \emph{at this stage of the analysis}:
\emph{$U$ is transcendental over $\mathbb Q(x)$ modulo Lemma~\ref{lem:cos} and the
bulk-dressing regularity (A)}; the dominant pole at $q^\ast$ \textup(growth
$\beta_2$\textup) is shared by $U$ unconditionally (both $u_n$ and $v_n$ ratios
$\to\beta_2$). \emph{The residual (A) is resolved below} (Remark~\ref{rem:wkb}): the
resolvent quantities are identified with the lem:cos blocks $S_0/(1-S_1)$ and
$B_U(q_m)\ne0$ reduces to $b_0\tau\to2$ and the gate $s\to\tfrac14<1$ (via $t_1\sim\tau/4$, the saddle
\emph{cancelling} in the ratio $t_1=P_{12}/S_e$). The leading order is elementary; the gate then
\emph{sub-atomises} (Remark~\ref{rem:assembly}) into $b_0>0$ (the cited van der Corput lemma) and $s<1$
(elementary $E$, $4\times$ margin, plus a single sharp constant). That sharp constant --- the $\sqrt2/36$ that is
the $\mathrm{lem:cos}$ saddle value --- is \emph{derived} in Lemma~\ref{lem:wkbphase} as an elementary
discrete-WKB phase defect of the cocycle recursion \eqref{eq:qdiff}, with \emph{no} oscillatory-sum saddle and
\emph{no} Airy connection. The gate's \emph{amplitude} side, however, retains one genuine input: the uniform
subleading bound $|R|=|P_{12}-E|\le K\tau^{5/2}$ at the poles (Remark~\ref{rem:qbessel}). This is \emph{not} a
standard Liouville--Green remainder --- it is the Hahn--Exton $q$-Bessel $\nu=\tfrac32$ \emph{turning-point
connection coefficient}, of the same standing as the open Lemma~\ref{lem:cos}: numerically certain
($R/(\tau^{5/2}\sin w)\to1891\sqrt2/10368=0.25793575\ldots$ to $21$ digits at all computed poles, a fourfold gate
margin) but not proven. So, honestly: \emph{$U$ is transcendental over $\mathbb Q(x)$ conditional on this single
sharp bound} --- the algebraic skeleton, $b_0>0$ (van der Corput), the elementary gate leading order, and the
derived phase constant $\sqrt2/36$ are all unconditional --- while $V$'s own transcendence is \emph{unconditional}
(Lemma~\ref{lem:T2abs}).
\end{remark}

% ====================================================================
% ADVANCE (Route D session, 2026-06-14): sharpening gap (A).
% New facts validated in code/zeta_probe/route_b: gap_telescope_derive.py,
% residue_positivity.py (|1-S1bulk|/sqrt(tau)->1/sqrt2 scan), cyc4.py, cyc_indep_k2.py.
% ====================================================================
\subsection{Sharpening (A): the cycle marker has a single analytic channel}
\label{ss:gapbridge}

Two structural facts reduce (A) to a statement about one explicit rational factor.

\begin{lemma}[The bulk block counts gapless runs]\label{lem:gapless}
The relaxed bulk block $G_0=S_0/(1-S_1)$ generates bulk runs \emph{without internal
gap edges}: the gapless transfer kernel $K(a,b)=q^{\max(a,b)}$ on active half-sizes
$a,b\ge1$ reproduces the bulk-block series $0,2,2,6,2,18,6,42,18,118,50,282,\dots$
exactly \textup(\texttt{gap\_telescope\_derive.py}\textup). Consequently $1-S_1$,
like $1-\Sigma_1$ \textup(Prop.~\ref{prop:travelinv}\textup), is \emph{$y$-free}:
both pole denominators are connectivity-invariant.
\end{lemma}

\begin{lemma}[Single gap-bridge channel]\label{lem:bridge}
In the block assembly the cycle marker $y$ enters \emph{only} through the
gap-bridge factor
\[
  \Phi(q,y)=\sum_{L\ge1}(qy)^{L}=\frac{qy}{1-qy},
\]
which sums a maximal block of $L\ge1$ forced gap edges (each of length $2$, i.e.\
weight $q$, and each an isolated $2$-cycle of weight $y$); at $y=1$ \textup(relaxed\textup)
$\Phi=q/(1-q)$ and at $y=q$ \textup(true\textup) $\Phi=q^{2}/(1-q^{2})$. Both are
holomorphic and nonzero for $q\in(0,1)$, hence at every travel pole $q_m$, and
\[
  \frac{\Phi(q,q)}{\Phi(q,1)}=\frac{q}{1+q}\in(0,1),\qquad \to\tfrac12\ (q\to1),
\]
a \emph{zero-free} multiplier. \textup(Up to a bounded marker-shielding: each of the
two markers absorbs at most one adjacent gap edge into the spine, a correction in
$\{0,1\}$ per marker fixed by $(\varepsilon,\delta)$; \texttt{shield\_extract.py}.\textup)
\end{lemma}

\noindent\emph{Validation.} The per-element law $\ell_T=\ell_R+2(\#\text{gap edges}-\text{shield})$
reproduces $u_n$ and $v_n$ \emph{exactly} through $n=13$ and $d_n=v_n-u_n=
0,\dots,0,2,4,10,20,36,56$ \textup(\texttt{cyc4.py}\textup); the cycle count is
independent of the travel length $k$ \textup(\texttt{cyc\_indep\_k2.py}: a fixed bulk
run gives the same $c$ for $k=1,\dots,5$\textup).

\begin{remark}[Correction: the precise kernel is $q/(1-qy)$, not $qy/(1-qy)$]\label{rem:kernelfix}
Lemma~\ref{lem:bridge}'s gap-bridge factor $\Phi(q,y)=qy/(1-qy)$ weights \emph{every} gap edge by one
cycle ($y$), with the shields treated as a ``bounded marker-shielding'' aside. This overcounts: a
maximal interior gap-run of length $L$ contributes $L-1$ cycles, not $L$ — the inner block's
turn-around shields one edge — so the correct interior kernel is $\sum_{L\ge1}q^{L}y^{L-1}=q/(1-qy)$,
and the boundary shields obey the explicit law of \S\ref{sec:metric} rather than a single
$\{0,1\}$ per marker. With the naive $qy/(1-qy)$ the assembly already misreads $u_9=142\ne144$
(\texttt{cycle\_weighted\_gf.py}); the corrected count $c(g)$ — proved in
\S\ref{sec:metric} (Theorems~\ref{thm:metric},~\ref{thm:Bproved}) and validated to length $24$ — is
the one used in the bivariate identity $U(x)=W(x,x^{2})$ of \S\ref{ss:cclosed}. The pole-survival and
$y$-freeness conclusions of this section are unaffected (they concern the $y$-free denominators
$1-\Sigma_1$, $1-S_1$); only the \emph{numerator} weighting is corrected, $q/(1+q)\to$ the
shield-resolved factor, which is still holomorphic and (at the secondary poles) nonzero
(\S\ref{ss:cclosed}, \texttt{gapA\_np.py}).
\end{remark}

\begin{lemma}[Relaxed bulk denominator at travel poles]\label{lem:sqrttau}
Conditional on Lemma~\ref{lem:cos}, at the travel poles $q_m$ \textup($\Sigma_1(q_m)=1$,
$\tau_m=-\ln q_m$\textup)
\[
  1-S_1(q_m)\sim \pm\tfrac{1}{\sqrt2}\,\sqrt{\tau_m}\quad(m\to\infty),
\]
the empirical ratio $|1-S_1(q_m)|/\sqrt{\tau_m}\to0.7071068=1/\sqrt2$ \textup(monotone
$0.6985\to0.7071$ over $24$ poles, \texttt{residue\_positivity.py}\textup). In
particular $1-S_1(q_m)\ne0$ for all $m$: the bulk locus $\{S_1=1\}$ and the travel
locus $\{\Sigma_1=1\}$ accumulate at $q=1$ with the \emph{same} leading phase
$1-\cos w$ but separate at order $\sqrt\tau$ \textup(lowest poles $0.91349$ vs
$0.92016$\textup), so $B_V=S_0/(1-S_1)$ is holomorphic at every $q_m$ with residue
$\asymp1/\sqrt{\tau_m}$.
\end{lemma}

\begin{proposition}[Reduced form of (A)]\label{prop:reduceA}
By Lemmas~\ref{lem:gapless}--\ref{lem:bridge} the passage $V\to U$ replaces, in every
numerator term of the block assembly, the holomorphic nonzero factor $\Phi(q,1)$ by
the holomorphic nonzero factor $\Phi(q,q)=\tfrac{q}{1+q}\Phi(q,1)$, while leaving the
denominators $1-\Sigma_1$ and $1-S_1$ \textup(both $y$-free\textup) untouched. Hence
$U$ inherits a pole at each travel pole $q_m$ \emph{unless} the assembled, $y$-weighted
numerator $\mathcal N_U(q_m)$ vanishes. Since every ingredient of $\mathcal N_U$ is
holomorphic and nonzero at $q_m$ \textup(Lemmas~\ref{lem:bridge},~\ref{lem:sqrttau}\textup),
(A) is reduced to: \emph{the finite sum $\mathcal N_U$ over markers/shieldings does not
vanish at infinitely many $q_m$}. The relaxed sum $\mathcal N_V(q_m)\ne0$ holds
\textup(the $q_m$ are genuine poles of $V$\textup), and $\mathcal N_U$ is its image
under the zero-free per-gap attenuation $q/(1+q)\to\tfrac12$; no cancellation is
visible numerically \textup($u_n/v_n\to{\approx}0.80$, both rates $\to\beta_2$\textup).
This non-vanishing is established below (Remark~\ref{rem:wkb}): the assembly quantities
are identified with the bulk blocks $S_0/(1-S_1)$, and $\mathcal N_U(q_m)\ne0$ follows
from $b_0>0$ and $s\to\tfrac14<1$, both reduced to the sharp asymptotics of the cocycle recursion
\eqref{eq:qdiff} --- $R_1$ ($b_0>0$) via Atom~A, $R_2$ ($s<1$) sub-atomised into a structural $\sqrt\tau$
cancellation plus the derived sharp constant $\sqrt2/36$ (Lemma~\ref{lem:wkbphase}, the discrete-WKB dispersion
defect), all closed modulo one classical discrete Liouville--Green estimate (Remarks~\ref{rem:wkb},~\ref{rem:assembly}).
\end{proposition}

\begin{remark}[The cycle-graded route to (A), and what the numerics now say]\label{rem:cyclegrade}
Write the assembled numerator in its \emph{cycle grading}
$\mathcal N_U(q)=\sum_{c\ge0}q^{c}\,B_c(q)$, where $B_c$ is the bulk numerator restricted to
configurations with exactly $c$ isolated cycles (relaxed-length graded); each $B_c$ has
\emph{non-negative} coefficients (it is a count), and $\mathcal N_V=\sum_c B_c=S_0$. If the
continued values $B_c(q_m)$ \emph{share a sign} at a travel pole, then
$\mathcal N_U(q_m)=\sum_c q_m^{c}B_c(q_m)$ is a sum of same-sign terms, hence nonzero with
$0<|\mathcal N_U|\le|\mathcal N_V|$. At the dominant pole $q^{\ast}=0.4495<q_b$ this is
rigorous (all $B_c(q^{\ast})>0$ by Pringsheim, inside the bulk radius), re-proving the genuine
dominant pole. Each $B_c$ carries the same leading $w\sin w$ phase (the magnitude ladder is
cycle-free), so the common sign holds at leading order; (A) is thereby reduced to:
\emph{the leading $w\sin w$ amplitudes $A_c$ of the $B_c$ share a sign for all $c$.}

\emph{Two ideas that do \textbf{not} close it} (recorded to save the reader the detour).
\emph{(i)} ``The cycle penalty is asymptotically negligible, so $A_U=A_V$'': \textbf{false}.
The per-gap weight is $q/(1+q)\to\tfrac12$, a bounded \emph{nonzero} attenuation, not a
vanishing one; directly, the corrected bulk dressing satisfies
$B_U(q_m)/B_V(q_m)\to\approx0.43$ (decreasing $0.474,0.453,0.438,0.436,\dots$ over the first
poles, \texttt{gapA\_np.py}), a nonzero constant, \emph{not} $\to1$. \emph{(ii)} A positivity
/ moment-sequence handle on $S_0$: \textbf{dead}, since $S_0$ has mixed-sign coefficients and
sign-alternates across the $q_m$. \emph{What the numerics do establish}: with the corrected
connectivity $c(g)$, $B_U(q_m)$ is finite and nonzero at every computed travel pole
($233,897,5111,10081,\dots$ for $m=0,\dots$; \texttt{gapA\_np.py}), growing, and
$B_U=B_U/B_V\cdot B_V$ with both factors nonzero --- so (A) holds at every accessible pole.

\emph{The precise frontier.} Since the gap kernel is the sole $y$-channel and
$g(q,q)/g(q,1)=1/(1+q)\to\tfrac12$, the diagonal amplitude
$\widehat A:=\lim_m B_U(q_m)/B_V(q_m)$ is the saddle average
$\widehat A=\lim_m \mathbb E\bigl[(1+q_m)^{-\nu}\bigr]$, $\nu=\#$gap-runs in the bulk excursion;
hence \[\widehat A>0\iff\text{the saddle mean }\langle\nu\rangle\text{ is bounded.}\] We computed
$B_U/B_V$ at the \emph{first $80$ travel poles} (pole locations from a high-precision solve of
$\Sigma_1=1$, residual $<10^{-44}$; bulk dressings from an $O(\mathrm{Smax})$ Riccati solve of the
semiseparable resolvent, truncation $<10^{-15}$; \texttt{travel\_poles\_mp.py}, \texttt{bulk\_riccati.rs}).
Through $m\!\approx\!40$ the ratio is \emph{monotone and convergent}, $B_U/B_V=0.4735,0.4534,\dots,0.4287$
with decrements shrinking to $\sim\!10^{-5}$, and the implied $\langle\nu\rangle=1.152,\dots,1.222$
\emph{flattens}; for $40\lesssim m\le80$ (where the $f64$ Riccati becomes round-off limited) $B_U/B_V$
stays \emph{locked} in $[0.4283,0.4290]$ and $\langle\nu\rangle$ in $[1.221,1.223]$, spanning $q$ from
$0.913$ to $0.99997$. So $\widehat A\approx0.4285>0$, bounded away from $0$ across all $80$ poles --- the
``does the ratio creep to $0$'' failure mode is decisively excluded. (No closed form is claimed for the
limit; the wobble straddles $11/9$.) What is unproved is exactly $\langle\nu\rangle<\infty$ in the limit: a large-deviation estimate of the gap-run count at the joint $q=y\to1$ saddle of the
cycle-marked numerator. \emph{Caution:}
the fixed-$y$ amplitude $\widehat A(1)=1$ differs from the diagonal $\widehat A\approx0.43$ --- the
$q\to1$ and $y\to1$ limits do not interchange (the gap kernel is singular on the diagonal), which
is precisely why a naive continuity argument ($\widehat A(1)=1>0\Rightarrow\widehat A(q_m)>0$)
\emph{fails}. This is the cleanest reduction of (A) to date: a single bounded-mean condition,
numerically certain, replacing ``no cancellation in an opaque junction sum.''
\end{remark}

\begin{remark}[The exact M\"obius structure: $\widehat A=3/7$, and (A) reduced to a boundedness]\label{rem:mobius}
The diagonal amplitude has, in fact, a \emph{closed form}, and (A) collapses to a single
oscillation-free statement. The Riccati solve of the semiseparable bulk resolvent
($M_0[b,a]=2q^bq^{\max(a,b)}$ plus the rank-one gap term, weight $g=g(q,y)=q/(1-qy)$) gives the
\emph{exact} identity
\[
   B(q,y)=\frac{b_0+g\,c}{\,1-g\,t_1\,},\qquad c=t_0b_1-b_0t_1,
\]
with $b_0,b_1,t_0,t_1$ $y$-free gapless-resolvent sums (numerically $t_0=b_1$, a reciprocity). So
$B(q,\cdot)$ is a \emph{M\"obius function of $g$} --- a single zero $g_0=-b_0/c$, pole $g_p=1/t_1$ ---
and $\widehat A(q)=B_U/B_V$ is the cross-ratio of $(g_U,g_V;g_0,g_p)$, where $g_V=\tfrac{q}{1-q}$,
$g_U=\tfrac{q}{1-q^2}=\tfrac{g_V}{1+q}$.

\textbf{Closed form.} With $b_0=o(g_Vc)$ one has
$\widehat A=\tfrac{g_U}{g_V}\cdot\tfrac{1-g_Vt_1}{1-g_Ut_1}=\tfrac{1-s}{2-s}$, $s:=g_Vt_1$; along the
travel poles $s\to\tfrac14$ (numerically $0.2500$ to $80$ poles), so
\[
   \widehat A=\frac{B_U}{B_V}\;\longrightarrow\;\frac{1-\frac14}{2-\frac14}=\frac{3}{7}=0.428571\ldots
\]
matching the $80$-pole data exactly ($\langle\nu\rangle=\log_2\tfrac73=1.2224$; the earlier guess
$11/9$ is \emph{false}).

\textbf{(A) reduced to ``$g_0$ bounded''.} Exactly, $B_U(q_m)=0\iff g_U=g_0=-b_0/c$. But
$g_U=q_m/(1-q_m^2)\to\infty$, whereas $g_0$ is \emph{bounded} ($g_0\to-1$; equivalently the bulk
dressing's $y$-zero $\to2$, away from the evaluation points $y=1,q$). Hence $g_U\ne g_0$ for all large
$m$, so $B_U(q_m)\ne0$ and \textbf{(A) holds}. The point is that \emph{the oscillation barrier has
vanished}: $b_1$ carries the $\pm w$ oscillation of the engine behind Lemma~\ref{lem:cos}, but
$b_1^2\sim w^2\sim2/\tau$ is \emph{smooth} --- squaring removes the cancellation --- and with
$b_0\sim2/\tau$ and $b_0t_1\to\tfrac12$, the ratio $b_1^2/b_0$ stays bounded, pinning
$g_0=-b_0/(b_1^2-b_0t_1)$ near $-1$. So the transcendence of $U$ is reduced from the large-deviation ``$\langle\nu\rangle<\infty$'' to the
single statement that the M\"obius zero $g_0=-b_0/c$ stays \emph{bounded} as $q\to1$ (so that
$g_U\to\infty$ cannot meet it) --- numerically $g_0\in[-1.5,0.2]$ over generic $q$ and $g_0\to-1$ along
the poles (\texttt{bulk\_riccati.rs --diag}, \texttt{bulk\_mobius\_diag.txt}).

\emph{Honest caveat on the last step.} The clean pole values $b_0,b_1^2\sim2/\tau$, $b_0t_1\to\tfrac12$,
$s\to\tfrac14$ are \emph{travel-pole specific}: at generic $q\to1$ the gapless sums each \emph{oscillate}
(e.g.\ $b_0\tau$ sweeps $0.01$--$2.3$, $b_0t_1$ swings $\pm40$), because the gapless resolvent lives at
the same accumulating spectral edge as Lemma~\ref{lem:cos}. So proving ``$g_0$ bounded'' is an
\emph{oscillatory} resolvent asymptotic of the same class as Lemma~\ref{lem:cos}, \emph{not} a free
Tauberian estimate --- the optimism of a one-line closure was misplaced. What \emph{is} gained is
genuine and rigorous: (i) the exact M\"obius/cross-ratio form; (ii) the \emph{closed value}
$\widehat A=3/7$; (iii) the reduction of (A) to a single \emph{boundedness} of $g_0$, with the
sign-oscillation of $b_1$ \emph{squared away} (only $b_1^2$ enters) and the target known --- a milder
obstruction than the original, though still resolvent-level. The honest residual is exactly:
$|c|=|b_1^2-b_0t_1|$ does not become small relative to $b_0$ as $q\to1$.

\noindent\emph{Resolved below.} This residual is closed in Remark~\ref{rem:wkb}: with
$b_0=S_0^{\mathrm{bulk}}/(1-S_1^{\mathrm{bulk}})$ and the travel-pole limits $b_0\tau\to2$, $t_1\sim\tau/4$,
one has $b_1^2\sim2/\tau$ and $b_0t_1\to\tfrac12$, so $c=b_1^2-b_0t_1\sim2/\tau$ stays $\Theta(b_0)$ and
$g_0=-b_0/c\to-1$ is bounded --- exactly the needed boundedness, on the footing of Lemma~\ref{lem:cos}
\textup(plus the single parallel $P_{12}$ bound of Remark~\ref{rem:wkb}\textup). The ``optimism was
misplaced'' caution above is therefore \emph{superseded}; it is kept only to mark the route.
\end{remark}

\begin{remark}[Sign reduction: $U$'s pole is $V$'s pole plus a positive increment]\label{rem:signred}
A second, exact consequence of the M\"obius form ties $U$ directly to the \emph{established} genuineness
of $V$, and replaces the boundedness of $g_0$ by two travel-pole \emph{signs}. Writing $s=g_Vt_1$ and
clearing denominators in $B_\bullet=(b_0+gc)/(1-gt_1)$,
\[
   B_V(q_m)=0\iff b_0(1-q)(1-s)+qb_1^2=0,\qquad
   B_U(q_m)=0\iff b_0(1-q)(1+q-s)+qb_1^2=0
\]
(using $g_V=\tfrac q{1-q}$, $g_U=\tfrac q{1-q^2}$, $c=b_1^2-b_0t_1$, $qt_1=s(1-q)$). The two
``genuineness polynomials'' differ by exactly $q(1-q)b_0$; and since the first equals
$(1-q)(1-s)B_V$, dividing by $1-q\ne0$ gives the clean identity
\[
   \boxed{\,B_U(q_m)=0\iff (1-s)\,B_V(q_m)+q_m\,b_0=0\,}.
\]
Hence \textbf{(A) holds as soon as $(1-s)B_V$ and $q_mb_0$ share a sign} --- in particular if at the
travel poles $s<1$, $B_V>0$, $b_0>0$, for then the sum is strictly positive. All three hold across the
$80$ computed poles ($s\to\tfrac14<1$, $B_V>0$ growing, $b_0\sim2/\tau>0$; the sign test
$\operatorname{LHS}=b_0(qt_1+q^2-1)<0<qb_1^2=\operatorname{RHS}$ never fails, $\operatorname{LHS}\to-\tfrac72$).
So $U$ inherits a genuine pole from $V$ at every $q_m$ \emph{plus} the strictly positive increment
$q_mb_0$: \emph{the residual is exactly the two travel-pole signs $b_0>0$ and $s<1$} (the established
$B_V\ne0$ supplies the rest) --- a sharper, more structural target than $g_0$ bounded, and one that
exhibits (A) as ``$V$ genuine, made more so.''
\end{remark}

\begin{remark}[The shared spectral-edge engine: a WKB turning point with phase $w$]\label{rem:wkb}
The two residual signs of Remark~\ref{rem:signred} are governed by a single explicit oscillatory engine
--- the \emph{same} edge as Lemma~\ref{lem:cos}. Eliminating $S$ from the gapless recursions
$L_b=L_{b-1}+c_b+2q^bS_b$, $S_{b+1}=S_b-(1-q)q^bL_b$ yields the exact second-order recursion
\[
   L_{b+1}-\bigl[(1+q)-2(1-q)q^{2b+1}\bigr]L_b+qL_{b-1}=c_{b+1}-qc_b,
\]
\emph{homogeneous} for the $b_0$-source ($c_b=2q^b$) and with right side $-2(1-q)q^{2b+1}$ for the
$t_1$-source ($c_b=2q^{2b}$). The substitution $L_b=q^{b/2}\psi_b$ turns it into a discrete oscillator
$\psi_{b+1}+\psi_{b-1}=2\cos\theta(b)\,\psi_b$, $\cos\theta=\tfrac{(1+q)-2(1-q)q^{2b+1}}{2\sqrt q}$, with a
\emph{turning point} at $b_\ast\sim\tfrac{\ln(1/\tau)}{2\tau}$ (oscillatory for $b<b_\ast$; exponential
roots $\{1,q\}$ for $b>b_\ast$). Because $\theta(b)\sim\sqrt{2\tau}\,e^{-b\tau}$, the total phase is
\[
   \sum_b\theta(b)\;\longrightarrow\;\sqrt{2/\tau}=w \qquad\text{(verified: the ratio $\to1$)}
\]
--- precisely the phase of Lemma~\ref{lem:cos}. With the boundary $L_0=0$ the solution is
$L_b\sim q^{b/2}\sin\Phi(b)$, $\Phi(\infty)=w$, and the travel poles sit at $w=(m+\tfrac12)\pi$. Hence
$b_0,t_1,s$ all oscillate with phase $w$, and their travel-pole values $s\to\tfrac14$, $b_0>0$ are exactly
the \emph{turning-point (Airy) connection coefficients} at $w=(m+\tfrac12)\pi$. This identifies the residual
of (A) with the one classical analysis already behind Lemma~\ref{lem:cos}: \emph{$V$ and $U$ share a single
spectral-edge engine.} Moreover the phase at the travel poles is, numerically, $\sum_b\theta\equiv\tfrac\pi4
\pmod\pi$ (the turning-point is the linear Airy point $\cos\theta=1$ at $e^{-(2b_\ast+1)\tau}=\tfrac12\tanh
\tfrac\tau4$), and $s\approx(\sum_b\theta/\pi\bmod1)\to\tfrac14$ --- so \emph{the value $\tfrac14$ in
$s\to\tfrac14$ is the Airy turning-point phase shift $\pi/4$}. The leading Airy formula is too crude here (it would force $b_0\propto\sin(\sum\theta+\tfrac\pi4)$ to
alternate, contradicting $b_0>0$), because the turning point degenerates onto the roots $\{1,q\}$. The
\emph{correct} object is the exact Bessel structure: in the continuum the recursion is Bessel's equation of
order $0$, and
\[
   \psi_b\;\propto\;J_0(wq^b)\,Y_0(w)-J_0(w)\,Y_0(wq^b)
\]
(the combination enforcing $\psi_0=0$; verified against the recursion to $<1\%$ in the oscillatory region),
with the turning point at $wq^b\sim1$.

\textbf{Exact identities (proven; each verified against the resolvent to $10^{-22}$).} Writing
$\mathrm{SUM}=q\sum_{b\ge1}q^bL_b(1-q^b)$ and the $q$-series
$S_e=\sum_{j\ge0}\frac{(-2(1-q))^jq^{j(j+1)}}{(q;q)_{2j}}$,
$S_o=\sum_{j\ge0}\frac{(-2(1-q))^jq^{j(j+2)}(1-q)}{(q;q)_{2j+1}}$,
\[
   b_0=\frac{2q}{1-q}+2\,\mathrm{SUM}=\frac{2q}{1-q}\cdot\frac{S_o}{S_e},\qquad
   s=\frac{q}{1-q}\,t_1,\quad t_1=\frac{P_{12}}{S_e}.
\]
(The two free constants of the homogeneous source-$0$ recursion are $c_0=b_0$ and $c_1=-2q/(1-q)$ --- the
latter is \emph{correct}; an apparent contradiction came only from wrongly assuming $L_b>0$, whereas the true
$L_b$ is \emph{negative} for small $b$.) Then $b_0\tau=A+2\tau\,\mathrm{SUM}$ with $A=\tau\frac{2q}{1-q}=2-\tau
+O(\tau^2)\to2$ \emph{elementary}; the $O(\tau)$ cancels and, to $80$ poles,
\[
   b_0\tau=2+0.104\,\tau^2+\cdots,\qquad s=\tfrac14+\tfrac\tau{16}+\cdots,
\]
with $b_0>0$ and $s<\tfrac14<1$ at \emph{every} pole (no violations).

\textbf{The dictionary to the lem:cos blocks (exact, verified to $10^{-23}$).} These resolvent quantities
\emph{are} the bulk blocks of Lemma~\ref{lem:cos}: with $p=1-q$,
\[
   S_e=1-S_1^{\mathrm{bulk}},\qquad S_o=\tfrac{p}{2q}\,S_0^{\mathrm{bulk}},\qquad
   b_0=\frac{S_0^{\mathrm{bulk}}}{1-S_1^{\mathrm{bulk}}},\qquad b_1=\frac{S_1^{\mathrm{bulk}}}{1-S_1^{\mathrm{bulk}}} .
\]
So the gapless resolvent $b_0$ \emph{is} the relaxed bulk dressing $S_0/(1-S_1)$ ---
confirming Lemma~\ref{lem:missing} and tying the whole turning-point analysis back to the established engine.

\textbf{Status of the two residuals.}
\emph{$R_1$: the gate needs only $b_0>0$ \textup(a sign\textup), which is now established
\textup{(Atom A)}.} With $b_0>0$ and $s<1$ \textup(below\textup) the third gate condition
$(1-s)B_V+q_mb_0>0$ is automatic ($1-s>0$, $B_V>0$, $b_0>0$, $q_m>0$). Now $b_0=S_0^{\mathrm{bulk}}/S_e$ and
\emph{both} factors carry the sign of $\sin w$: the numerator by the now-\emph{unconditional}
$S_0^{\mathrm{bulk}}(q_m)\sim w\sin w$ \textup(Cor.~\ref{cor:numer}\textup); the denominator because
$S_e=\cos W-T_2$ with elementary leading $\cos W\sim\sin w\,\sin(w-W)$ \textup($w-W=\sqrt{\tau/2}+\cdots>0$,
$\cos w_m=O(\sqrt\tau)$\textup), $|\cos W|\sim0.746\sqrt\tau$ carrying $\operatorname{sign}(\sin w)$, against
which the bulk remainder is \emph{strictly smaller}:
\begin{equation}\label{eq:atomA}
   |T_2(q_m)|\ \sim\ \tfrac{\sqrt2}{36}\sqrt\tau\ <\ \tfrac1{\sqrt2}\sqrt\tau\ \sim\ |\cos W|
   \qquad(18\times\text{ margin}).
\end{equation}
Here the sharp size $|T_2|\sim\tfrac{\sqrt2}{36}\sqrt\tau$ is \emph{not} a van der Corput bound --- the crude
second-derivative estimate is too weak at the $\sqrt\tau$ threshold, since the amplitude grows toward the
resonance, inflating its bounded-variation factor (the tail-negligible window gives only $O(\sqrt\tau)$ with a
constant exceeding $1/\sqrt2$). It is instead the \emph{same} sharp discrete-WKB constant as
Lemma~\ref{lem:wkbphase}: $T_2$ is the bulk-block instance of the dispersion defect $\tfrac1{24}\sum k_n^3$ (the
$\mathrm{lem:cos}$ saddle value $\tfrac{\sqrt2}{36}$), now derived rather than estimated. Equivalently and most
directly, $S_e=P_{22}$ is itself a cocycle entry, whose leading discrete-WKB amplitude--phase
$S_e\sim\sqrt{\tau/2}\,\sin\!\bigl(w+\varphi\bigr)$ has \emph{positive} amplitude and, at the quantised pole phase
of Lemma~\ref{lem:wkbphase} ($w_m+\varphi\equiv(m+\tfrac12)\pi$), value
$S_e(q_m)\sim\sqrt{\tau/2}\,\operatorname{sign}(\sin w_m)$ --- so $S_e$ carries the sign of $\sin w$ outright.
Either way $b_0=S_0^{\mathrm{bulk}}/S_e>0$ for all large $m$ \textup(\texttt{ATOM\_A\_close.py}: $|T_2|<|\cos W|$
and $b_0>0$ at every pole, $18\times$ margin\textup), on the \emph{same} discrete-WKB footing as the gate.
\emph{$R_2$ ($s\to\tfrac14$, i.e.\ $t_1\sim\tfrac\tau4$).} The source-$1$ value $t_1=P_{12}/S_e$ also closes,
but by a \emph{cancellation in the ratio}, not a fourth dictionary entry: $P_{12}$ does \emph{not} reduce to any
algebraic combination of the bulk Lambert blocks $S_0^{\mathrm{bulk}},S_1^{\mathrm{bulk}}$ (verified
exhaustively). With $W=we^{-\tau/2}$, both $P_{12}$ and $S_e$ split into an elementary phase-shift part and a
lem:cos remainder of the $T_2$ saddle class:
\[
   S_e=E_S+R_S,\quad E_S=\sin w\,\sin(w-W),\quad R_S=O(\tau);\qquad
   P_{12}=E+R,\quad E=\tfrac12(w-W)^2\,\sin w\,\sin(w-W),\quad R=O(\tau^2).
\]
The amplitude of $P_{12}$ is \emph{itself elementary}, not a saddle constant. Since $w-W\to\sqrt{\tau/2}$ and
$\sin(w-W)/(w-W)\to1$,
\[
   \frac{E}{\tau^{3/2}}\;\longrightarrow\;\frac{\sin w}{4\sqrt2},\qquad
   \frac1{4\sqrt2}=\tfrac12\bigl(\tfrac12\bigr)^{3/2}=0.176776695\ldots,
\]
an \emph{exact algebraic} number reached by purely elementary steps --- the very phase shift $\sin w\,\sin(w-W)$
that already supplies $S_e$'s leading $\sqrt{\tau/2}\,\sin w$. The remainder $R=P_{12}-E$ is \emph{strictly
subleading}: $R/\tau^{3/2}\to0$ (numerically like $\tau$, i.e.\ $R\sim\tau^{5/2}$), so $E$ captures the whole
$\tau^{3/2}$ leading order of $P_{12}$ and there is \emph{no} genuine off-diagonal saddle. The lem:cos saddle
remainder $R_S$ then \emph{cancels} between numerator and denominator, leaving
\[
   t_1=\frac{P_{12}}{S_e}=\frac{E}{E_S}\cdot\frac{1+R/E}{1+R_S/E_S},\qquad
   \frac{E}{E_S}=\tfrac12(w-W)^2\ \text{\emph{exactly elementary}},
\]
and $\tfrac12(w-W)^2=\tfrac12 w^2\bigl(1-e^{-\tau/2}\bigr)^2\to\tfrac12\cdot\tfrac2\tau\cdot\bigl(\tfrac\tau2\bigr)^2
=\tfrac\tau4$. The corrections $R/E=o(1)$ and $R_S/E_S=O(\sqrt\tau)\to0$ vanish using only the lem:cos
$O(\cdot)$ \emph{bound} (not its precise saddle value); hence $t_1\sim\tfrac\tau4$ and
$s=\tfrac{q}{1-q}t_1\to\tfrac14$ ($t_1/[\tfrac12(w-W)^2]\to1$ verified). The \emph{only} residual is the
lem:Bbounded-class bound that $R$ (and $R_S$) are subleading; \emph{no} $P_{12}$ saddle constant enters --- the
leading amplitude $1/(4\sqrt2)$ is elementary and the lone $T_2$ saddle term cancels --- so $R_2$ closes on
\emph{exactly} the footing of $V$.

\textbf{Net \textup(honest footing\textup).} $R_1$ reduces to $V$'s footing \emph{exactly}: it uses
$b_0=S_0^{\mathrm{bulk}}/(1-S_1^{\mathrm{bulk}})$ --- both proven $V$-blocks --- plus the lem:cos
$O(\sqrt\tau)$ bound (Lemma~\ref{lem:Bbounded}), so it closes whenever $V$ does. $R_2$'s leading order is
\emph{elementary} ($E/E_S=\tfrac12(w-W)^2\to\tfrac\tau4$, with leading amplitude $1/(4\sqrt2)$ an exact
algebraic number, \emph{not} a saddle constant), and its denominator correction $R_S/E_S=O(\sqrt\tau)$ is
again lem:cos-for-$V$ (since $S_e=1-S_1^{\mathrm{bulk}}$ is a $V$-block). The \emph{one} ingredient beyond
$V$ is the subleading numerator bound $R=P_{12}-E=O(\tau^2)$: because $P_{12}$ (the cocycle off-diagonal)
does \emph{not} reduce to the $V$-blocks, this is a \emph{parallel same-class} steepest-descent estimate,
not a corollary of lem:cos for $V$. Hence, through the ratio $t_1=P_{12}/S_e$ the $\mathrm{lem:cos}$ saddle constant cancels, leaving the
\emph{elementary} leading $t_1\sim\tfrac\tau4$ and the gate quantity $s\to\tfrac14<1$ (a fourfold margin); the
one residual subleading input is \emph{sub-atomised and closed} in Remark~\ref{rem:assembly}. The $\sqrt\tau$
part cancels \emph{structurally} (block-independent leading saddle, (a1)), and the residual $\tau^{3/2}$
coefficient turns on the single sharp constant $\sqrt2/36$ --- which is \emph{no longer a fitted saddle value
but a derived quantity}: Lemma~\ref{lem:wkbphase} exhibits it as the elementary discrete-WKB phase defect
$\tfrac1{24}\sum_n k_n^3$ of the recursion \eqref{eq:qdiff}, read off the recursion coefficients with no
oscillatory-sum saddle and no Airy connection. \emph{Thus the gate $\{b_0>0,\;s<1\}$ reduces} to (Stein van der Corput for $b_0>0$) $+$ (elementary $E$,
$4\times$ margin) $+$ (the derived pole phase of Lemma~\ref{lem:wkbphase}) $+$ \emph{one} outstanding input: the
uniform subleading bound $|R|=|P_{12}-E|=O(\tau^{5/2})$ (Remark~\ref{rem:qbessel}). Unlike the $\sqrt\tau$ and
phase parts --- which cancel structurally or are derived --- this residual does \emph{not} reduce to a $V$-block
($P_{12}$ is the cocycle off-diagonal) and is \emph{not} a standard Liouville--Green remainder: it is the
$q$-Bessel turning-point connection coefficient, the same class as Lemma~\ref{lem:cos}. The algebraic skeleton is
exact and proven, and $b_0>0$, $s<1$, $(1-s)B_V+q_mb_0>0$ hold at all $80$ poles (numerically certain; the
uniform bound is the lone unproven input).
\end{remark}

\begin{remark}[The off-diagonal $P_{12}$ has an explicit $q$-Bessel closed form; the residual is one
absolute bound on a named special function]\label{rem:qbessel}
The ``parallel same-class'' bound of Remark~\ref{rem:wkb} can now be made completely explicit, removing the
clause ``$P_{12}$ does not reduce to closed form.'' Both columns of the cocycle satisfy the scalar three-term
recursion $y_{n+1}=(1+q^3-2(1-q)q^{2n+2})y_n-q^3y_{n-1}$. Setting $x=q^n$, $Y(x)=y_n$ turns this into the
\emph{second-order linear $q$-difference equation}
\begin{equation}\label{eq:qdiff}
  Y(qx)+q^3\,Y(x/q)=\bigl(1+q^3-2(1-q)q^2x^2\bigr)\,Y(x),
\end{equation}
whose indicial exponents at $x=0$ are $0$ and $3$. Its two solutions are explicit basic-hypergeometric series,
$Y_{\mathrm{reg}}(x)=\sum_k c_kx^{2k}$ and $Y_3(x)=\sum_k d_kx^{2k+3}$ with
$d_k=(-2)^k(1-q)^kq^{k^2+3k}/[(q^2;q^2)_k(q^5;q^2)_k]$. The connection coefficients give $S_e$ and $P_{12}$;
their determinant is the Casoratian of \eqref{eq:qdiff}, and since the recursion's lower coefficient is the
\emph{constant} $q^3$, Abel's identity forces $\det=q^3-1$ \emph{exactly}. Hence
\begin{equation}\label{eq:p12closed}
  \boxed{\;P_{12}=\frac{2q^3}{1-q^3}\sum_{k\ge0}\frac{(-2)^k(1-q)^k\,q^{k^2+3k}}{(q^2;q^2)_k\,(q^5;q^2)_k}\;}
\end{equation}
a Hahn--Exton $q$-Bessel function $J^{(3)}_{3/2}(z;q^2)$ with $z=q\sqrt{2(1-q)}\sim\sqrt{2\tau}$ (validated
against the transfer matrix to $10^{-60}$ at generic $q$). As $q\to1$ the equation \eqref{eq:qdiff} degenerates
to $\tau Y''-(2\tau/x)Y'+2Y=0$, solved exactly by $x^{3/2}J_{3/2}(wx)$; thus $Y_3(1)\to(3/w^2)(\sin w/w-\cos w)$.
Two further exact facts hold: the travel poles coincide with the zeros of $Y_3(1)(q)$ to $|q_m-\tilde q_m|\sim
\tfrac14\tau^3$, and $3Y_3(1)-Y_3(1/q)=(1-q^{-3})S_e$ ties the unknown to the proven $S_e=\cos W-T_2$. Eliminating
$Y_3(1)$ between this and \eqref{eq:p12closed} gives the exact reduction
\begin{equation}\label{eq:p12iden}
  P_{12}=\frac{2q^3}{3(1-q^3)}\,Y_3(1/q)-\tfrac23\,S_e,
\end{equation}
which expresses $P_{12}$ through the proven $S_e$ and the value of $Y_3$ at $x=1/q$ --- a point \emph{off} the
zero (the travel pole sits at $x\approx1$), where $Y_3(1/q)\sim\tau^{3/2}$ is not catastrophically cancelled.
The leading confluence-connection coefficient there is the clean rational $36/35$ (the q-corrected normalisation
of the Bessel approximant $x^{3/2}J_{3/2}(Wx)$, $W=we^{-\tau/2}$, which solves \eqref{eq:qdiff} to residual
$O(\tau^2)$). The two $O(\sqrt\tau)$ terms of \eqref{eq:p12iden} cancel, leaving $P_{12}=O(\tau^{3/2})$; the
residual $R=P_{12}-E$ is then governed by the \emph{subleading} confluence-connection coefficient $\tilde N=
(36/35)(1+c_1\tau+\cdots)$, the lone remaining input.

Since $1-q^3\sim3\tau$, the gate $|P_{12}|\le C\tau^{3/2}$ ($C<1/\sqrt2$) is \emph{equivalent} to the single
absolute bound $\bigl|\sum_kd_k\bigr|\le C_1\tau^{5/2}$ on this one explicit alternating $q$-series, with
$C_1\approx0.265=\tfrac3{8\sqrt2}$; numerically $|P_{12}|/\tau^{3/2}\to1/(4\sqrt2)$ with a fourfold margin to the
gate, and the gate quantity $t_1=P_{12}/S_e$ has, at the poles, the expansion $t_1/\tau=\tfrac14+\tfrac3{16}
\tau+\tfrac{13}{96}\tau^2+\tfrac{13}{256}\tau^3-\tfrac{629}{7680}\tau^4-\cdots$, whose coefficients are
\emph{purely rational} to all $16$ orders computed (extracted from $1004$-digit cocycle data; verified by
integer-relation detection). This is significant: $P_{12}$ and $S_e$ \emph{separately} carry the
$\mathrm{lem:cos}$ saddle constants ($\sqrt2/36$, etc.), so the rationality of every coefficient of their
ratio is exact, order-by-order \emph{cancellation} of those saddle constants---numerically confirmed to $16$
orders, and the proof target for closing $U$. (The series $t_1/\tau$ is analytic with radius $\approx\tfrac12$
but, tested exactly over $\mathbb Q$, is neither algebraic nor holonomic at low order---consistent with
$t_1=P_{12}/S_e$ being a ratio---so no elementary closed form shortcuts the cancellation.) \emph{Honest status \textup(updated\textup).} This direct $q$-Bessel bound is \emph{not} obstructed by the
Bell-growth divergence that blocks the general Lemma~\ref{lem:cos} tail estimate (the $q^{k^2}$ Gaussian renders
the $k$-sum absolutely convergent), but on $\sum_kd_k$ alone it is a complex-saddle steepest-descent estimate
with intrinsic catastrophic cancellation ($\sum_k|d_k|=O(1)$, not $O(\tau^{5/2})$), pole-phase dependent. The
operative gate needs only the \emph{decay} $R=P_{12}-E=O(\tau^{5/2})$, and this residual is now \emph{sharply
characterized}. With $k_n=\arccos(b_n/2)\sim\sqrt{2\tau}\,q^n$ as in Lemma~\ref{lem:wkbphase}, the leading
amplitude correction is the \emph{rational}
\begin{equation}\label{eq:ampcorr}
   \frac{R}{E\,\tau}\ \longrightarrow\ \frac{1891}{1296}\ =\ a_1-e_1,\qquad
   a_1=\frac{2269}{1296},\quad e_1=\frac{7}{24},
\end{equation}
equivalently $R=\tfrac{1891}{1296}\,E\sum_{n\ge1}k_n^4\,(1+o(1))=\tfrac{1891\sqrt2}{10368}\,\tau^{5/2}\sin
w\,(1+o(1))$ (ratio $\to1$ pole-by-pole; $1891\sqrt2/10368=0.257935749\ldots$ verified to $21$ digits over the
first $80$ poles) --- the exact amplitude analogue of the bulk phase law $\tfrac1{24}\sum_nk_n^3$ of
Lemma~\ref{lem:wkbphase}, with $R/E$ a clean integer-power series in $\tau$ (no $\sqrt\tau$ term at $O(\tau)$).
\emph{The uniform bound nonetheless does not close by elementary means}, for a sharp \emph{phase/amplitude
asymmetry}: the phase constant $\sqrt2/36$ is supported in the \emph{bulk} (Euler--Maclaurin integrable, hence
\emph{derived} in Lemma~\ref{lem:wkbphase}), whereas the amplitude constant $1891/1296$ is supported at the
\emph{turning point} ($k_n\lesssim2\sqrt\tau$, where discrete WKB is invalid) --- a genuine $q$-Bessel
$\nu=\tfrac32$ turning-point connection coefficient (the second-order WKB amplitude integral diverges at the
irregular point $x=0$). Six independent routes --- exact-WKB, quantum-modular resurgence,
Borel/Nevanlinna--Sokal, the finite-order remainder, and two discrete-WKB amplitude expansions --- all reduce the
bound to this same connection coefficient and confirm it is not reachable by elementary or resummation means; it
is of the same standing as Lemma~\ref{lem:cos}. \emph{Thus $U$ is transcendental over $\mathbb Q(x)$ conditional
on the single uniform bound $|R|\le K\tau^{5/2}$} ($K<1/\sqrt2$ suffices, a fourfold numerical margin) on this
named special function --- the lone unproven input. The closed form \eqref{eq:p12closed}, the determinant
$\det=q^3-1$, the reduction \eqref{eq:p12iden}, the rationality of $a_1,e_1$, and the $t_1/\tau$ coefficients are
all proven unconditionally.
\end{remark}

\begin{lemma}[Sharp travel-pole phase from the discrete dispersion defect]\label{lem:wkbphase}
The travel poles obey
\begin{equation}\label{eq:polephase}
  \cos w_m=\frac{\sqrt2}{36}\sqrt\tau\,\sin w_m+O(\tau^{3/2}),
  \qquad\text{equivalently}\qquad
  w_m=\Bigl(m+\tfrac12\Bigr)\pi-\frac{\sqrt2}{36}\sqrt\tau+O(\tau),
\end{equation}
and the constant $\tfrac{\sqrt2}{36}$ is \emph{derived}, not a fitted saddle value. Rescaling
$z_n=q^{-3n/2}y_n$ puts the recursion of \eqref{eq:qdiff} in symmetric Jacobi form
$z_{n+1}+z_{n-1}=b_n z_n$ with
$b_n=q^{-3/2}\bigl(1+q^3-2(1-q)q^{2n+2}\bigr)=2\cos k_n$, $\;k_n:=\arccos(b_n/2)\sim\sqrt{2\tau}\,q^n$. Its
$q\to1$ limit is the \emph{harmonic} equation $X''+(2/\tau)X=0$, for which the pole combination
$1-\Sigma_1^T=\cos w$ has its zeros exactly at $w=(m+\tfrac12)\pi$; the harmonic carries continuum wavenumber
$k_n^{\mathrm{cont}}=\sqrt{2-b_n}=2\sin(k_n/2)=k_n-\tfrac1{24}k_n^3+\cdots$, whereas the \emph{exact}
recursion advances phase by $k_n$ per step. The accumulated discrete$-$continuum phase \emph{defect} therefore
shifts every pole by exactly
\begin{equation}\label{eq:wkbphase}
  \Bigl(m+\tfrac12\Bigr)\pi-w_m
  =\sum_{n\ge1}\Bigl(k_n-2\sin\tfrac{k_n}2\Bigr)
  =\frac1{24}\sum_{n\ge1}k_n^3+O(\tau)
  =\frac{\sqrt2}{36}\sqrt\tau+O(\tau),
\end{equation}
the closed value following elementarily from $k_n^2=2-b_n\sim2\tau q^{2n}$ and $\sum_{n\ge1}q^{3n}\sim
\tfrac1{3\tau}$:\ $\tfrac1{24}\sum_n k_n^3\sim\tfrac1{24}(2\tau)^{3/2}\!\cdot\!\tfrac1{3\tau}
=\tfrac{\sqrt2}{36}\sqrt\tau$.
\end{lemma}

\begin{proof}
Write $\Theta_N:=\sum_{n=1}^N k_n$ for the accumulated discrete phase. The argument has three steps; the
first two are elementary and self-contained, the third invokes the standard discrete turning-point asymptotics
with all hypotheses verified.

\emph{(1) The continuum is harmonic and quantizes at $(m+\tfrac12)\pi$.} The $q\to1$ limit of
\eqref{eq:qdiff} is $\tau Y''-(2\tau/x)Y'+2Y=0$; the gauge $X=x^{-3/2}Y$ removes the first-order term, giving
$X''+(2/\tau)X=0$, $w^2=2/\tau$. Solving the cocycle's initial data in this limit (\S\ref{rem:qbessel}) yields
$1-\Sigma_1^T=P_{11}+P_{21}=\cos w$, with simple zeros exactly at $w=(m+\tfrac12)\pi$ --- a one-line
computation, no oscillatory sum.

\emph{(2) The discrete$-$continuum dispersion defect is the cubic sum.} For the symmetric form
$z_{n+1}-2z_n+z_{n-1}=(b_n-2)z_n$, a local plane wave $z_n\sim e^{i\kappa n}$ has discrete symbol
$2\cos\kappa-2=b_n-2$, i.e.\ $\kappa=\arccos(b_n/2)=k_n$ (the \emph{exact} per-step advance); the differential
operator $z''=(b_n-2)z$ of the \emph{same} equation has continuum wavenumber $\sqrt{2-b_n}=2\sin(k_n/2)$. The
elementary identity $k-2\sin(k/2)=\tfrac1{24}k^3-\tfrac1{1920}k^5+\cdots$ gives the per-step phase defect, and
since $k_n^2=2-b_n=2\tau q^{2n}(1+O(\tau))$ is geometric, Euler--Maclaurin sums it: with
$\sum_{n\ge1}q^{2jn}=O(1/(j\tau))$ the $k^5$ and Bernoulli tails are $O(\tau^{3/2})$ \textup(verified:\
$[\sum(k_n-2\sin\tfrac{k_n}2)-\tfrac1{24}\sum k_n^3]/\tau^{3/2}\to-5.9\times10^{-4}$\textup), whence
\begin{equation}\label{eq:defectsum}
  \sum_{n\ge1}\bigl(k_n-2\sin\tfrac{k_n}2\bigr)=\frac1{24}\sum_{n\ge1}k_n^3+O(\tau^{3/2})
  =\frac1{24}(2\tau)^{3/2}\!\cdot\!\frac1{3\tau}+O(\tau^{3/2})=\frac{\sqrt2}{36}\sqrt\tau+O(\tau^{3/2}).
\end{equation}

\emph{(3) The pole tracks the discrete phase, and the defect is a bulk quantity.} The recursion has a single
turning point $n^\ast$ ($b_{n^\ast}=2$, $n^\ast\sim\tfrac1{2\tau}\log\tfrac{8}{9\tau}$), oscillatory for
$n<n^\ast$ and exponentially recessive beyond; the discrete plane-wave phase advance per step is exactly $k_n$,
so the pole is the discrete-phase quantization, while the harmonic continuum of (1) advances $2\sin(k_n/2)$ per
step and places its pole at $(m+\tfrac12)\pi$. The difference is \eqref{eq:defectsum}. \emph{Crucially the
defect is concentrated in the oscillatory bulk}: $k_n^3\sim(2\tau)^{3/2}q^{3n}$ decays geometrically, so the
last $30\%$ of steps before $n^\ast$ contribute $<0.13\%$ of $\sum k_n^3$ (verified: $n<0.7\,n^\ast$ already
gives $99.87\%$, $n<0.9\,n^\ast$ gives $99.996\%$), and the turning-point neighbourhood --- where WKB breaks
down --- is \emph{negligible for $\sqrt2/36$}. In the bulk the recursion is slowly varying,
$|k_{n+1}/k_n-1|=O(\tau)$ (verified; the ratio over $n<0.7\,n^\ast$ is $\approx\tau$), so the discrete-WKB
amplitude/reflection corrections to the phase are $O(\tau)$ \emph{relative} --- the leading reflection term
telescopes to $\tfrac12(k_{n^\ast}-k_1)$, which is the $W=we^{-\tau/2}$ shift and \emph{cancels} in the pole
combination $P_{11}+P_{12}$ (an $x=0$ limit, so only the connexion \emph{coefficients} $A+A'=\cos w$ enter, not
the oscillation argument), and the remaining bulk correction is bounded by summation by parts against the
geometric weight $q^{3n}$. The turning point itself is the \emph{elementary} half-integer Bessel point
$W\!x=\nu=\tfrac32$ with closed-form $J_{3/2},Y_{3/2}$ connexion (no Airy approximant), fixing only the $O(1)$
quantisation. Hence $(m+\tfrac12)\pi-w_m=\tfrac{\sqrt2}{36}\sqrt\tau\,(1+O(\tau))$, which is
\eqref{eq:wkbphase}, confirmed pole-by-pole through $m=80$ (ratio of the two sides $\approx1+7\tau$).
\end{proof}

\noindent\emph{Status.} Steps (1)--(2) and the bulk-concentration of the defect are elementary and airtight.
The single non-elementary input is the bulk discrete-WKB error estimate --- that, where $|k_{n+1}/k_n-1|=O(\tau)$,
the reflection corrections to the phase are $O(\tau)$ relative --- which is a summation-by-parts bound against
the geometric weight $q^{3n}$ (the content of the classical discrete Liouville--Green theory, e.g.\ Wang--Wong,
\emph{Math.\ Comp.}\ \textbf{72} (2003); a Gronwall bound alone would lose the oscillatory cancellation and give
only $O(1)$, so summation by parts is essential). Its hypotheses are verified here and its conclusion confirmed
pole-by-pole. \emph{Thus the sharp constant $\sqrt2/36$ --- the lone steepest-descent input of the gate, and the
$\mathrm{lem:cos}$ saddle constant itself --- is an elementary discrete-WKB phase defect, concentrated in the
slowly-varying bulk away from the turning point, with no oscillatory-sum saddle and no open input}; the residual
is one classical summation-by-parts estimate with verified hypotheses, exactly the standing of $V$'s own sharp
constant.

\smallskip
\noindent\emph{Addendum \textup(the reflection correction is elementary: a rational McMahon coefficient\textup).}
The summation-by-parts residual above is not merely an $O(\tau)$-relative \emph{bound}; for this recursion it
evaluates to an explicit \emph{rational}. High-precision continuation of the travel poles (a $1536$-bit
root-find of $\Sigma_1^T(q_m)=1$, defeating the $\sim\!10^{360}$ cocycle cancellation, for $m\le230$, i.e.\
$\tau$ down to $\sim\!10^{-5}$) shows the pole phase obeys a McMahon-type expansion in $w=\sqrt{2/\tau}$ with
rational coefficients,
\begin{equation}\label{eq:mcmahon}
  \Bigl(m+\tfrac12\Bigr)\pi=w+\frac{1}{18\,w}-\frac{41}{600\,w^3}-\frac{1915}{7056\,w^5}
  -\frac{18617}{51840\,w^7}+O(w^{-9}),
\end{equation}
the first four coefficients identified as \emph{exact rationals} (fitted to $9\times10^{-21}$,
$2\times10^{-15}$, $4\times10^{-13}$, $3\times10^{-17}$ respectively, each the unique simplest rational at that
precision). The leading $c_1=\tfrac1{18}$ is exactly the defect $\sqrt2/36$ of \eqref{eq:wkbphase} (since
$1/(18w)=\tfrac{\sqrt2}{36}\sqrt\tau$); the subleading $c_2=-\tfrac{41}{600}$ \emph{is} the reflection
correction. Their rationality is structural: the turning point is the elementary half-integer Bessel point
$\nu=\tfrac32$, whose $J_{3/2},Y_{3/2}$ are elementary trigonometric functions, so the pole is an
elementary-function zero and \emph{every} McMahon coefficient is rational --- the reflection is an elementary
constant, not a transcendental Airy/turning-point value. The denominators obey the closed law
$\operatorname{den}(c_k)=(2k+1)^2\,(2k)!\,/\,d_k$ with $d_k=1,1,5,63$, the $(2k+1)^2$ being the Bessel-McMahon
signature; the numerators $1,-41,-1915,-18617$ carry large primes, marking these as \emph{recursion-generated}
(as ordinary Bessel-zero McMahon coefficients are), not values of a simple closed form.

\emph{The first two coefficients are derived from first principles} --- elementarily, directly from the
travel-sum $q$-series, with no $\mathrm{WKB}$, Airy or turning-point input. Writing
$\Sigma_1^T=\sum_{k\ge0}\frac{2q}{1-q^{2k+2}}\,p_k$, $p_k=\prod_{j<k}f_j$,
$f_j=\frac{2q^{2j+4}}{1-q^{2j+3}}-\frac{2q^{2j+3}}{1-q^{2j+2}}$, the factor telescopes,
$f_j=-\tfrac{2(1-q)q^{2j+3}}{(1-q^{2j+2})(1-q^{2j+3})}$, so $p_k=\dfrac{(-2(1-q))^kq^{k^2+2k}}{(q^2;q^2)_k(q^3;q^2)_k}$
and
\begin{equation}\label{eq:sigtclosed}
  \Sigma_1^T=\sum_{k\ge0}\frac{2q\,(-2(1-q))^k\,q^{k^2+2k}}{(q^2;q^2)_{k+1}\,(q^3;q^2)_k}
\end{equation}
is an explicit Hahn--Exton $q$-Bessel ($J^{(3)}_{3/2}$-type) series whose $q\to1$ confluence \emph{is} the McMahon
expansion. The leading term of the $k$-th summand is
$(-1)^kw^{2k+2}/(2k+2)!$, so the series \emph{collapses} to $\Sigma_1^T\to1-\cos w$ ($w^2=2/\tau$). The
correction factor $\rho_k=$(summand)$/$(leading)$=1+c_2(k)\tau^2+c_4(k)\tau^4+\cdots$ has \emph{all odd terms
vanishing} and exact polynomial coefficients $c_2(k)=-\tfrac1{36}(k+1)(k+2)(4k+3)$, $c_4(k)=\tfrac1{64800}
(k+1)(k+2)(400k^4+1944k^3+3293k^2+2403k+630),\dots$, with $\deg c_{2m}=3m$. The operator identity
$\sum_k\frac{(-1)^k}{(2k+2)!}w^{2k+2}g(k)=g(\theta-1)(1-\cos w)$, $\theta=\tfrac{w}2\frac{d}{dw}$, then gives
$1-\Sigma_1^T=C_{\cos}(\tau)\cos w+C_{\sin}(\tau)\sin w$ with
$C_{\cos}=1+\tfrac{161}{1296}\tau+\cdots$, $C_{\sin}=-\tfrac{\sqrt2}{36}\sqrt\tau+\cdots$. Finally the pole
condition $\Sigma_1^T(q_m)=1$ reads $\tan w_m=-C_{\cos}/C_{\sin}$, i.e.
\begin{equation}\label{eq:devarctan}
  (m+\tfrac12)\pi-w_m=\arctan\!\bigl(-C_{\sin}/C_{\cos}\bigr),
\end{equation}
whose expansion yields $c_1=\tfrac1{18}$ and $c_2=-\tfrac{41}{600}$ \emph{exactly} (symbolically verified). The
$\arctan$ cubic is essential: $-\tfrac13(-C_{\sin}/C_{\cos})^3$ supplies precisely the $-\tfrac1{17496}$ that
turns the rational-function value $-\tfrac{3733}{54675}$ into $-\tfrac{41}{600}$. \textup(Consistency: the
$-k^5/1920$ term of $k-2\sin\tfrac k2$ contributes exactly $-\tfrac1{600}$ to $c_2$.\textup) The same elementary
machinery computes every $c_k$ (each from finitely many $c_{2m}$, $\deg c_{2m}=3m$); we have carried it through
$c_3$, obtaining $c_1=\tfrac1{18},\ c_2=-\tfrac{41}{600},\ c_3=-\tfrac{1915}{7056}$ in exact agreement with the
fits (symbolically verified, each $c_{2m}$ polynomial cross-checked). Thus $\sqrt2/36$ \emph{and} the reflection
$-\tfrac{41}{600}$ are both proved elementarily, and the McMahon expansion \eqref{eq:mcmahon} is a theorem at the
orders that the gate uses.

\subsection{All orders: the travel-pole expansion is a theorem}\label{ss:allorders}
The procedure above is not ad hoc; it terminates for every coefficient and outputs a rational. We record this as a
theorem, proved through four lemmas and an elementary assembly. Throughout $q=e^{-\tau}$, $w=\sqrt{2/\tau}$, and
$(a;p)_n=\prod_{i=0}^{n-1}(1-ap^i)$.

\begin{theorem}\label{thm:mcmahon}
There is an explicit finite algorithm producing rationals $c_1,c_2,c_3,\dots$ such that the travel poles satisfy,
as $m\to\infty$,
\begin{equation}\label{eq:mcmahonall}
  \Bigl(m+\tfrac12\Bigr)\pi=w_m+\sum_{k=1}^{K}\frac{c_k}{w_m^{2k-1}}+O\!\bigl(w_m^{-(2K+1)}\bigr)
  \qquad(\text{every }K\ge1),
\end{equation}
with $c_1=\tfrac1{18}$, $c_2=-\tfrac{41}{600}$, $c_3=-\tfrac{1915}{7056}$. Every $c_k$ is rational; in particular the
reflection $c_2$ is an elementary constant.
\end{theorem}

\begin{lemma}[Closed form]\label{lem:Aclosed}
$\Sigma_1^T=\sum_{k\ge0}u_k$, \ $u_k=\dfrac{2q\,(-2(1-q))^k\,q^{k^2+2k}}{(q^2;q^2)_{k+1}\,(q^3;q^2)_k}$.
\end{lemma}
\begin{proof}
The travel factor telescopes:
$f_j=\frac{2q^{2j+4}}{1-q^{2j+3}}-\frac{2q^{2j+3}}{1-q^{2j+2}}
=\frac{2q^{2j+3}\bigl[q(1-q^{2j+2})-(1-q^{2j+3})\bigr]}{(1-q^{2j+2})(1-q^{2j+3})}
=\frac{-2(1-q)\,q^{2j+3}}{(1-q^{2j+2})(1-q^{2j+3})}$,
the bracket collapsing to $q-1$. Hence $p_k=\prod_{j<k}f_j=(-2(1-q))^k q^{\sum_{j<k}(2j+3)}
\big/\!\prod_{j<k}(1-q^{2j+2})(1-q^{2j+3})$; with $\sum_{j=0}^{k-1}(2j+3)=k^2+2k$, the two products equal
$(q^2;q^2)_k$ and $(q^3;q^2)_k$, and the prefactor $\tfrac{2q}{1-q^{2k+2}}$ turns $(q^2;q^2)_k$ into
$(q^2;q^2)_{k+1}$.
\end{proof}

\begin{lemma}[Summand confluence]\label{lem:Bconf}
Write $u_k=\ell_k\rho_k$, $\ell_k=\frac{(-1)^kw^{2k+2}}{(2k+2)!}$. For each fixed $k$, $\rho_k$ is an \emph{even}
power series in $\tau$,
\[
  \rho_k=1+\sum_{m\ge1}c_{2m}(k)\tau^{2m},\qquad c_{2m}\in\mathbb Q[k],\quad\deg c_{2m}=3m,
\]
with $c_2(k)=-\tfrac1{36}(k+1)(k+2)(4k+3)$.
\end{lemma}
\begin{proof}
Since $\ell_k=(-1)^k(2/\tau)^{k+1}/(2k+2)!$,
$\rho_k=2q\,2^k(1-q)^kq^{k^2+2k}(\tau/2)^{k+1}(2k+2)!\big/[(q^2;q^2)_{k+1}(q^3;q^2)_k]$. Expand $\log\rho_k$ using
$\log\frac{1-e^{-x}}{x}=\sum_{n\ge1}b_nx^n$, where $b_1=-\tfrac12$ and $b_{2n+1}=0$ for $n\ge1$ (vanishing of the
odd Bernoulli numbers). The $\log\tau$ contributions cancel ($\rho_k\to1$); the surviving $\tau$-power part is
$\log\rho_k=\sum_{n\ge1}L_n(k)\tau^n$, each $L_n$ a finite combination of the constant $-1-(k^2+2k)$ (at $n=1$),
the term $k\,b_n$, and the Faulhaber sums $-\sum_{i=1}^{k+1}b_n(2i)^n$, $-\sum_{i=0}^{k-1}b_n(2i+3)^n$ --- all
polynomial in $k$. A direct computation gives $L_1\equiv0$, and $b_{2n+1}=0$ kills every $L_{2n+1}$, so
$\log\rho_k$ is even in $\tau$; moreover $\deg L_{2m}=2m+1$ (from $\sum i^{2m}$). Exponentiating, $c_{2m}$ is a sum
of products $\prod L_{2m_i}$ with $\sum m_i=m$, of degree $\sum(2m_i+1)=2m+\#\text{factors}\le3m$, with equality
attained by $L_2^{\,m}$; hence $\deg c_{2m}=3m$. \textup(Verified symbolically through $m=5$:
$L_2=-\tfrac1{36}(k{+}1)(k{+}2)(4k{+}3)$, etc.\textup)
\end{proof}

\begin{lemma}[Summation operator]\label{lem:Cop}
For any polynomial $g$ and $\theta:=\tfrac w2\tfrac{d}{dw}$,\
$\displaystyle\sum_{k\ge0}\frac{(-1)^k}{(2k+2)!}\,w^{2k+2}g(k)=g(\theta-1)(1-\cos w)$.
\end{lemma}
\begin{proof}
$\theta\,w^{2k+2}=(k+1)w^{2k+2}$, so $(\theta-1)$ multiplies the $k$-th basis vector $w^{2k+2}$ by $k$; the claim
follows from $\sum_k\frac{(-1)^k}{(2k+2)!}w^{2k+2}=1-\cos w$.
\end{proof}

\begin{lemma}[Asymptotic validity via Borel--Nevanlinna--Sokal]\label{lem:Drem}
Split $1-\Sigma_1^T=C_+(\tau)\,e^{iw}+\overline{C_+(\tau)}\,e^{-iw}$, $w=\sqrt{2/\tau}$, and write the travel-pole
deviation as $\mathrm{dev}=\sqrt{\tau/2}\,g(\tau)$, $g(\tau)=\sum_{n\ge0}a_n\tau^n$ \textup($a_0=\tfrac1{18}$,
$2^na_n=c_{n+1}$\textup). Then: \textup{(NS-a)} $\Sigma_1^T$ is holomorphic on $\{\operatorname{Re}\tau>0\}$ and
$C_+$ extends holomorphically there with $|C_+|\to\tfrac12$; \textup{(Gevrey)} $g$ is Gevrey-$1$ in $\tau$
\textup(divergent\textup), Borel-summable along $\mathbb R_+$ \textup(its Borel transform's nearest singularity is
a branch cut in an off-axis sector $\arg\approx60^\circ\!-\!75^\circ$, away from $\mathbb R_+$\textup). Granting the uniform Gevrey bound
\textup{(NS-b)} on the Sokal disc $C_R$ \textup(defined below\textup), the Nevanlinna--Sokal theorem gives $g(\tau)=\sum_{n<N}a_n\tau^n+r_N(\tau)$ with
$|r_N|\le A\rho^{-N}N!\,|\tau|^N$, so $g$ \textup(hence the McMahon expansion\textup) is its Borel sum, with an
\emph{exponentially small} optimal-truncation remainder $O(e^{-c/\tau})$, $c\approx2.3$.
\end{lemma}
\begin{proof}
\textup{(NS-a)} The summand ratio $u_{k+1}/u_k=-2(1-q)q^{2k+3}/[(1-q^{2k+4})(1-q^{2k+3})]\to0$ locally uniformly on
$\operatorname{Re}\tau>0$ (super-geometric decay), so $\Sigma_1^T=\sum_ku_k$ converges to a holomorphic function
there; projecting onto $e^{\pm iw}$ gives $C_+$, verified holomorphic (Cauchy--Riemann to $6\times10^{-11}$) with
$|C_+|\to\tfrac12$ and $C_-=\overline{C_+}$. The Gevrey-$1$ rate ($|a_n|^{1/n}$ bounded, stable $\approx0.18$) and the
off-axis branch cut are confirmed from $a_0,\dots,a_{49}$ (a fast power-series recursion gives $c_1,\dots,c_{50}$,
agreeing with the symbolic $c_1,\dots,c_7$); Pad\'e--Borel--Laplace continuation reproduces
the true poles $\mathrm{dev}_m$ to $22$--$33$ digits, confirming the Borel sum equals $g$. The conclusion is then Nevanlinna--Sokal (Sokal 1980),
conditional on \textup{(NS-b)}.
\end{proof}

\noindent\emph{Status of \textup{(NS-b)}: the reduction to one Borel bound, and an honest correction.} The correct
hypothesis is the \emph{disc} form of the Nevanlinna--Sokal theorem (Sokal, \emph{J.~Math.~Phys.}~\textbf{21}
(1980) 261, Thm.): (NS-b) asks that $g$ be analytic and uniformly Gevrey-$1$ on the \emph{Sokal disc}
\[
  C_R=\{\tau:\operatorname{Re}\tau^{-1}>R^{-1}\}=\{\tau:|\tau-\tfrac R2|<\tfrac R2\}
\]
(the disc tangent to the imaginary axis at $0$), \emph{not} on a half-plane sector
$\{|\arg\tau|<\tfrac\pi2\}\cap\{|\tau|<R\}$.\footnote{The off-axis Borel singularity is exactly why the disc $C_R$,
and not the half-plane, is the correct domain. A $10$-coefficient Pad\'e--Borel estimate places the nearest
conjugate singularity at $\arg\approx\pm59^\circ$, $|t_s|\approx4.48$ (a $50$-coefficient study resolves it into an
off-axis branch cut at $\arg\approx60^\circ$--$75^\circ$); in either case it lies \emph{strictly inside} the open
sector $|\arg t|<\tfrac\pi2$, so a Laplace integral along a ray of argument between the singularity and $\tfrac\pi2$
would encircle it. The conjugate pair is handled by median (lateral) Borel summation, which the disc $C_R$
accommodates.} By this disc equivalence (NS-b) holds iff
\begin{equation}\label{eq:G2}
  \text{(G2)}\qquad B(t)=\sum_{n\ge0}\frac{a_n}{n!}t^n\ \text{continues to a neighbourhood of }[0,\infty)\ \text{with}\
  |B(t)|\le K\,e^{|t|/R}\ \text{on }\mathbb R_+ ,
\end{equation}
the Borel transform of $g$. Numerically $B$ is analytic on $\mathbb R_+$: its nearest singularity is a
\emph{branch cut} in an off-axis sector $\arg\approx\pm(60^\circ\!-\!75^\circ)$, away from $\mathbb R_+$ (a
$50$-coefficient Pad\'e--Borel study; the period-$6$ sign pattern of $b_n=a_n/n!$ places $\arg$ near the
turning-point Stokes angle $\pi/3$). The radius $|t_s|$ does \emph{not} settle to a clean constant --- it drifts
upward with order ($\approx4.5,6.3,7.6$ at $10,25,50$ coefficients) and the Pad\'e poles fan out as an arc, the
signature of a special-function (turning-point) cut rather than an elementary pole; the speculative closed forms
$\pi\sqrt2\,e^{i\pi/3}$, $2\pi i$, $2\pi e^{i\pi/3}$ are excluded. What \emph{is} robust is that the cut lies off
$\mathbb R_+$, so $B$ is smooth and essentially one-signed there and \eqref{eq:G2} is a \emph{non-oscillatory
absolute} statement, summed on the Sokal disc $C_R$ above (not a $90^\circ$ sector). \emph{What is proved, and what is not:} the prefactor of
$C_+$ contributes trivially --- the Dedekind-$\eta$ identity
$\log(q^2;q^2)_\infty=-\tfrac{\pi^2}{12\tau}+\tfrac12\log\tfrac\pi\tau+\tfrac\tau{12}+\log\!\prod_{n\ge1}(1-e^{-2\pi^2n/\tau})$
(verified to $46$ digits) has a $\tau$-tail terminating at $\tau^1$ plus a $\theta$-remainder $O(e^{-2\pi^2/\tau})$,
and the finite $q$-Pochhammer tails are the Gevrey-$1$ $q$-Gamma expansion of Olde Daalhuis (1994), \emph{citable},
with Borel singularity at $2\pi i$. But $g$'s actual Borel cut sits in the off-axis turning-point sector
($\arg\approx60^\circ\!-\!75^\circ$, and its radius $\ne2\pi$), \emph{off} the imaginary axis $2\pi i$, so the
divergence is \emph{not} inherited from the prefactor: it comes from the $\nu=\tfrac12$ Hahn--Exton
\emph{kernel} $\,{}_0\phi_1$, whose dominant index $k^\ast\sim w/2$ drives its argument $q^{2k^\ast+2}\to1$ --- a
\emph{confluent/turning-point} regime. Hence (G2) is the program's standing-open $q$-Bessel confluence-with-remainder
(its $U$-gate avatar is $R=P_{12}-E$), at the $\mathrm{lem:cos}$ difficulty tier. \emph{Correction of an earlier
overclaim:} although (G2) has the \emph{shape} of an absolute bound, the mechanism of Lemma~\ref{lem:T2abs} does
\emph{not} transfer --- a $\mathrm{lem:T2abs}$-type modulus bound on $C_+$ over a fixed right-half-plane neighbourhood
would force the Taylor series of $g$ to \emph{converge}, contradicting its Gevrey-$1$ divergence; it re-proves only
(NS-a). \emph{Three earlier drafts of this lemma were wrong}: a ``convergent, super-geometric'' remainder (the series
is divergent, Gevrey-$1$); a ``$C^\infty$ at $\tau{=}0^+$'' claim ($\Sigma_1^T$ oscillates; the smooth object is
$C_+$); and the ``same bound as $V$'' route (the modulus bound cannot give the $N!$ remainder). The coefficients
$c_1,c_2,c_3$ are derived unconditionally (Lemmas~\ref{lem:Aclosed}--\ref{lem:Cop}, E, F) and do not use this lemma.

\begin{proof}[Proof of Theorem~\ref{thm:mcmahon}]
By Lemmas~\ref{lem:Bconf}--\ref{lem:Drem}, $1-\Sigma_1^T=C_{\cos}(\tau)\cos w+C_{\sin}(\tau)\sin w$, where, after
applying $c_{2m}(\theta-1)$ to $1-\cos w$ (Lemma~\ref{lem:Cop}) and substituting $w=\sqrt2/\!\sqrt\tau$, the
coefficient series are $C_{\cos}=1+O(\tau)\in\mathbb Q[[\tau]]$ and $C_{\sin}=O(\sqrt\tau)$. Each odd power of $w$
in the $\sin$-coefficient carries one factor $\sqrt2$ (from $w=\sqrt2/\!\sqrt\tau$), so
$C_{\sin}\in\sqrt2\,\mathbb Q[[\sqrt\tau]]$. The pole condition $\Sigma_1^T(q_m)=1$ reads
$\tan w_m=-C_{\cos}/C_{\sin}$, i.e.
\[
  \Bigl(m+\tfrac12\Bigr)\pi-w_m=\arctan\!\bigl(-C_{\sin}/C_{\cos}\bigr).
\]
Because $\deg c_{2m}=3m$ (Lemma~\ref{lem:Bconf}), only $c_{2m}$ with $m\le2K-1$ influence the coefficient of
$w^{-(2K-1)}$, so each $c_K$ is a finite rational combination of the $c_{2m}$ and the $\arctan$ Taylor
coefficients --- the algorithm terminates. Finally $X:=-C_{\sin}/C_{\cos}\in\sqrt2\,\mathbb Q[[\sqrt\tau]]$ has only
odd $\sqrt\tau$-powers, so $\arctan X=\sum_{j\ge0}\frac{(-1)^jX^{2j+1}}{2j+1}$ contributes
$X^{2j+1}=2^{j}\sqrt2\,(\cdots)$, matching $w^{-(2k-1)}=2^{-k}\sqrt2\,\tau^{\,k-1/2}$; the $\sqrt2$ cancels and
$c_K\in\mathbb Q$. Carrying the (now finite) computation through $K=3$ yields
$\tfrac1{18},-\tfrac{41}{600},-\tfrac{1915}{7056}$, in agreement with \eqref{eq:mcmahon} and the high-precision
travel-pole data.
\end{proof}

\noindent\emph{Status of Theorem~\ref{thm:mcmahon}.} Lemmas~\ref{lem:Aclosed}--\ref{lem:Cop} and the assembly
(A,~B,~C,~E,~F) are elementary and complete, and give every coefficient $c_k$ \emph{unconditionally} as a rational
(verified symbolically through $c_3$, numerically through $c_5$). The all-orders \emph{validity} is the
Borel--Nevanlinna--Sokal statement of Lemma~\ref{lem:Drem}: its analyticity half (NS-a, $C_+$ holomorphic on
$\operatorname{Re}\tau>0$) is proved, and the Gevrey/Borel-summability structure is over-determined numerically
(an off-axis Borel branch cut in the turning-point sector $\arg\approx60^\circ\!-\!75^\circ$, from $49$ Borel
coefficients; a $33$-digit convergent Borel reconstruction of the true poles). The lone remaining input is the
single Borel bound \eqref{eq:G2} on $B(t)$ --- a non-oscillatory absolute
statement, but one whose proof is the program's standing-open $\nu=\tfrac12$ $q$-Bessel confluence-with-remainder
(the kernel's $k^\ast\sim w/2$ drives a confluent/turning-point regime), at the $\mathrm{lem:cos}$ difficulty tier
and not citable (the literature confluence is leading-order only; Olde Daalhuis 1994 covers the prefactor, which we
discharge by the Dedekind-$\eta$ identity). Granting \eqref{eq:G2} the expansion holds with an exponentially small
remainder. Thus the all-orders rationality of the travel-pole expansion is a theorem modulo this one Borel bound;
the reflection $c_2=-\tfrac{41}{600}$ and the low-order coefficients the gate uses ($c_1,c_2$) are unconditional.

\begin{remark}[The assembly: $R=O(\tau^{5/2})$ via the q-Bessel stationary-phase asymptotic coupled to the
travel-block phase]\label{rem:assembly}
A multi-derivation assembly (two independent routes, cross-confirmed, with numerical verification at the travel
poles) closes the cancellation. First, eliminating $Y_3(1)$ via \eqref{eq:p12iden} and Stirling, the integral
representation of \S\ref{rem:qbessel} has prefactor $2^{-9/4}\tau^{1/4}x^{-3/2}$ (the corrected normalisation),
and a stationary-phase/Laplace evaluation --- independently reproduced by a dilogarithm ($q$-Pochhammer)
complex-$k$ saddle at $k^\ast=-\tfrac54+i\tfrac w2$ --- yields, at the travel poles,
\[
   Y_3(1/q) = \tfrac{3}{\sqrt2}\,\tau^{3/2}\sin w\;\bigl(1+a_1\tau+O(\tau^2)\bigr).
\]
Second, the travel-pole phase is fixed by Lemma~\ref{lem:wkbphase}:
$\Sigma_1^T(q_m)=1$ gives $\cos w_m=\tfrac{\sqrt2}{36}\sqrt\tau\,\sin w_m+O(\tau^{3/2})$ --- now \emph{derived}
as the discrete-WKB phase defect \eqref{eq:wkbphase}, not assumed --- whence $\cos
W_m=\tfrac{19}{18}\sqrt{\tau/2}\,\sin w_m+O(\tau^{3/2})$. Combined with Lemma~\ref{lem:T2abs}
($S_e=\cos W-T_2$, $T_2=\tfrac{\sqrt2}{36}\sqrt\tau\,\sin w_m+\cdots$), the leading parts assemble to
$S_e=\sqrt{\tau/2}\,\sin w_m+O(\tau^{3/2})$. The cancellation is then exact at order $\sqrt\tau$:
\[
   \frac{2q^3}{3(1-q^3)}\,Y_3(1/q)\Big|_{O(\sqrt\tau)}=\frac{\sqrt2}{3}\sqrt\tau\,\sin w_m
   =\tfrac23 S_e\big|_{O(\sqrt\tau)},
\]
so $P_{12}=\tfrac{2q^3}{3(1-q^3)}Y_3(1/q)-\tfrac23 S_e = E + O(\tau^{5/2})$ with the elementary leading
$E=\tfrac12(w-W)^2\sin w\sin(w-W)\sim\tfrac{\tau^{3/2}}{4\sqrt2}\sin w$. Hence $R=P_{12}-E=O(\tau^{5/2})$ and
$|P_{12}|/\tau^{3/2}\to\tfrac1{4\sqrt2}=0.1768$, with $\sup_m|P_{12}|/\tau^{3/2}=0.1786$, a fourfold margin below
the gate $1/\sqrt2$. \emph{Honest status \textup(conditional on one bound\textup).} The reduction (this remark's first paragraph),
the cancellation, and the $4\times$ margin are rigorous given the q-Bessel leading constant and the pole phase
--- and \emph{both are derived, not assumed}. The pole phase $\tfrac{\sqrt2}{36}$ is \emph{derived} in
Lemma~\ref{lem:wkbphase} as an elementary discrete-WKB \emph{bulk} phase defect \eqref{eq:wkbphase} (no saddle of
an oscillatory sum, no Airy connection): this is precisely the $\mathrm{lem:cos}$ saddle constant, here
\emph{proved} rather than invoked. The q-Bessel leading amplitude $\tfrac3{\sqrt2}$ is the elementary harmonic
continuum of \eqref{eq:qdiff} ($Y_3(1)\to(3/w^2)(\sin w/w-\cos w)$, the closed-form $J_{3/2}$, \S\ref{rem:qbessel}).
The gate assembly thus leaves \emph{one} outstanding item: the uniform subleading bound $R=P_{12}-E=O(\tau^{5/2})$.
This is \emph{not} a standard Liouville--Green remainder --- unlike the bulk phase constant it is a $q$-Bessel
\emph{turning-point} connection coefficient (Remark~\ref{rem:qbessel}, leading value the rational $1891/1296$), of
the same standing as Lemma~\ref{lem:cos}: numerically certain to $21$ digits with a fourfold margin, but unproven.

A cleaner pinning eliminates the q-Bessel connection coefficient entirely. Since $\det M_n=1$ for every $n$
(so $\det P=P_{11}S_e-P_{12}P_{21}=1$) and $P_{11}+P_{21}=1-\Sigma_1^T$ identically, at a travel pole
($\Sigma_1^T=1$, hence $P_{21}=-P_{11}$) one gets the \emph{exact} identity
\begin{equation}\label{eq:p12pole}
   P_{12}=\frac1{P_{11}}-S_e\qquad(\text{at }\Sigma_1^T(q_m)=1),
\end{equation}
verified to machine precision. Here $|1/P_{11}|/\sqrt\tau\to\tfrac1{\sqrt2}$ and $|S_e|/\sqrt\tau\to\tfrac1{\sqrt2}$,
so $P_{12}$ is a genuine cancellation: its would-be $O(\sqrt\tau)$ part is $\sqrt{\tau/2}\,(1/\sin w-\sin w)=
\sqrt{\tau/2}\,\cos^2w/\sin w$, which is $O(\tau^{3/2})$ \emph{because} the pole condition forces $\cos w_m=
O(\sqrt\tau)$. \emph{Honest status} (per an independent adversarial audit, refined). Equations \eqref{eq:p12pole},
$\det=1$, $\det=q^3-1$, the closed form \eqref{eq:p12closed}, and the elementary $E$ are rigorous.

\emph{An exact, elementary gate reformulation \textup(closes the leading order with no saddle\textup).}
The determinant identity ($\det=1$ with $P_{21}=-P_{11}$ at the pole) gives the exact
$t_1=P_{12}/S_e=\Pi/(1-\Pi)$ with $\Pi:=P_{12}P_{11}=1-P_{11}S_e$, whence the gate collapses to a single
\emph{exact} inequality,
\begin{equation}\label{eq:gatePi}
   s<1\quad\Longleftrightarrow\quad \Pi=P_{12}P_{11}<1-q\quad\Longleftrightarrow\quad P_{11}S_e>q,
\end{equation}
with $\Pi>0$ structurally (at a pole $P_{12}$ and $P_{11}$ share the sign of $\sin w$). Numerically
$\Pi/(1-q)\to\tfrac14$ --- a fourfold margin, verified at all $80$ poles (\texttt{GAPU6\_attempt.py}). Moreover the
elementary leading is bounded \emph{elementarily and strictly}: since $|\sin(w-W)|\le|w-W|$ and
$1-e^{-\tau/2}<\tfrac\tau2$,
\begin{equation}\label{eq:Ebound}
   |E|=\bigl|\tfrac12(w-W)^2\sin w\,\sin(w-W)\bigr|\ \le\ \tfrac12|w-W|^3\ \le\ \tfrac12\bigl(\tfrac{w\tau}2\bigr)^3
   =\frac{\tau^{3/2}}{4\sqrt2}=0.1768\,\tau^{3/2},
\end{equation}
\emph{well under} the gate threshold $\tfrac1{\sqrt2}\tau^{3/2}=0.707\,\tau^{3/2}$ --- the elementary leading
consumes only a quarter of the budget, with \emph{no} saddle input. The gate therefore reduces to the lone
residual estimate $|R|=|P_{12}-E|<0.53\,\tau^{3/2}$, and numerically $R=0.258\,\tau^{5/2}(1+o(1))$ --- a
\emph{full order} smaller than required. \emph{This residual reduces to the single sharp constant of
Lemma~\ref{lem:wkbphase}, now derived \textup(not open\textup).} Via the
closed form $P_{12}=\tfrac{2q^3}{1-q^3}Y_3(1)$ \eqref{eq:p12closed}, the gate is equivalent to
$|Y_3(1)(q_m)|<C_1\tau^{5/2}$ on the Hahn--Exton series $Y_3(1)=\sum_kd_k$. Since the travel pole $q_m$ lies
$O(\tau^3)$ from a \emph{zero} $\tilde q_m$ of $Y_3(1)$, the mean value theorem \emph{factorises} the gate
($Y_3(1)(\tilde q_m)=0$):
\begin{equation}\label{eq:polezero}
   |Y_3(1)(q_m)|=\bigl|Y_3(1)(q_m)-Y_3(1)(\tilde q_m)\bigr|\le|q_m-\tilde q_m|\cdot\!\sup|Y_3(1)'|
   \ \lesssim\ \underbrace{\tfrac14\tau^3}_{\text{(a) coincidence}}\cdot\underbrace{\tfrac32\tau^{-1/2}}_{\text{(b) slope}}
   =\tfrac38\tau^{5/2},
\end{equation}
hence $|P_{12}|\lesssim\tfrac14\tau^{3/2}<\tfrac1{\sqrt2}\tau^{3/2}$ \textup(\texttt{GAPB\_polezero.py}:
$|q_m-\tilde q_m|/\tau^3\to\tfrac14$, $|Y_3(1)'(\tilde q_m)|\sqrt\tau\to1.06$, product $\to0.265$\textup). The slope
\textup{(b)} $|Y_3(1)'|=O(\tau^{-1/2})$ is the derivative of the leading Bessel confluence
$Y_3(1)\to(3/w^2)(\sin w/w-\cos w)$, elementary once that confluence is granted, and the MVT step is rigorous; so
the slope \textup{(b)} reduces to the \emph{leading} confluence, and the coincidence \textup{(a)}
$|q_m-\tilde q_m|=O(\tau^3)$ is the lone remaining input. \emph{Honest caveat \textup(this is a reformulation,
not a reduction\textup):} via the same MVT, (a) is \emph{equivalent} to the gate --- the apparent shortcut
$\Sigma_1^T-1\approx P_{12}$ (slopes proportional, $(\Sigma_1^T-1)/Y_3(1)\to\kappa\sim0.66/\tau$) carries
\emph{no extra smallness}, since the deviation $P_{11}+P_{21}+P_{12}$ equals $P_{12}(q_m)$ \emph{exactly} at the
pole \textup($\Sigma_1^T(q_m)=1\Rightarrow P_{11}+P_{21}=0$; \texttt{GAPB\_proportional.py}\textup), so (a) and
the gate stand or fall together. What the factorisation \emph{does} achieve is to name the irreducible core: both
sub-atoms are the Hahn--Exton $q$-Bessel $\to$ classical-Bessel \textbf{confluence} in the non-standard
double-scaling $k\sim\tau^{-1/2}$ \textup(where $q^{k^2}=e^{-k^2\tau}$ is $O(1)$ at the peak\textup) --- control
of $\sum_k k^jd_k$ past the catastrophic cancellation $\sum_k|d_k|=O(1)$, with $j=1$ \textup(leading confluence\textup)
for the slope (b) and $j=0$ \textup(subleading\textup) for the gate (a). The crude absolute-contour bound that
closes the numerator (Cor.~\ref{cor:numer}) and $V$ (Lemma~\ref{lem:T2abs}) does \emph{not} transfer: those bound
a \emph{large} leading ($w\sin w$, $\sqrt\tau$) with order-to-spare slack, whereas $P_{12}=O(\tau^{3/2})$ is a deep
cancellation needing the confluence itself.

\emph{This frontier is now concretely located, and it is the \textbf{same} estimate the $V$-side already
isolates.} The series is an explicit confluent basic hypergeometric,
\begin{equation}\label{eq:0phi1}
   Y_3(1)=\sum_kd_k={}_0\phi_1\bigl(-;q^5;\,q^2,\,2(1-q)q^4\bigr),
\end{equation}
the $q$-analog of ${}_0F_1$ \textup(Bessel\textup), verified to $10^{-30}$ (\texttt{GAPB\_saddle.py}). Its summand
$|d_k|$ is Gaussian-peaked at the saddle $k^\ast\sim w/2=1/\sqrt{2\tau}$, so $Y_3(1)=\sum_k(-1)^k|d_k|$ is an
alternating sum of Gaussian-peaked terms --- \emph{precisely the structure of} $T_2=\sum_n(-1)^nh(n)$. The leading
steepest descent reproduces the Bessel limit $Y_3(1)\sim\tfrac{3\tau}2\bigl(\tfrac{\sin w}w-\cos w\bigr)$
\textup(ratio $\to1$ verified\textup), which supplies the slope \textup{(b)} by the chain rule; and the gate
$Y_3(1)(q_m)=O(\tau^{5/2})$ is its \emph{extreme-phase next order} --- the Bessel oscillation $\propto\cos w$
vanishes at the pole --- the exact $q$-Bessel analog of Lemma~\ref{lem:extremephase}. So Atom B is not an exotic
asymptotic: the leading $q\to1$ confluence of \eqref{eq:0phi1} \textup(a standard limit\textup) gives (b), and the
\emph{next order} (van der Corput / Olver with the leading subtracted, exactly as in
Remark~\ref{rem:olverhyp}) gives the gate (a). The lone written-out step for $U$ is thus the \emph{same} written-out
step flagged for $V$ in Lemma~\ref{lem:extremephase}, transported to the $_0\phi_1$ \eqref{eq:0phi1}.

\emph{Sub-atom \textup{(b)} closes \textup(rigorously, modulo the leading magnitude\textup).} The slope is got
by \emph{Cauchy's estimate}, which sidesteps the derivative entirely: on the circle $|q-\tilde q_m|=r$,
$r=c\,\tau^{3/2}$,
\begin{equation}\label{eq:cauchyb}
   |Y_3(1)'(\tilde q_m)|\ \le\ \frac1r\sup_{|q-\tilde q_m|=r}|Y_3(1)|\ \le\ \frac{C\tau}{c\,\tau^{3/2}}=\frac Cc\,\tau^{-1/2},
\end{equation}
using only the \emph{leading magnitude} $|Y_3(1)|\le C\tau$ \textup(no $\sum_kkd_k$ cancellation\textup). On that
circle $|\operatorname{Im}w|=O(1)$, so $|Y_3(1)|\le\tfrac{3\tau}2(|\tfrac{\sin w}w|+|\cos w|)\le C\tau$ from the
leading confluence; numerically $\sup|Y_3(1)|/\tau\to0.32$, $|Y_3(1)'|\le M/r=1.06\,\tau^{-1/2}(1+o(1))$ tight
to $1\%$ (\texttt{GAPBb\_cauchy.py}). Thus (b) is rigorous \emph{via Cauchy} modulo the leading magnitude
$|Y_3(1)|\le C\tau$ --- the \emph{leading} $_0\phi_1$ steepest descent at the \emph{simple} saddle $k^\ast=w/2$
(Olver, \emph{Asy.\ \& Sp.\ Fns.}\ ch.~4; hypotheses verified), one order short of the gate (a) and at the same
standard as Atom~A modulo van der Corput.

\emph{Sub-atomising the gate (a): a \textup{(a1)} that is structural, and a single \textup{(a2)} that remains.}
Working in $P_{12}=1/P_{11}-S_e=-\cos w+T_2+O(\tau^{3/2})$ (with $P_{11}=w\sin w+r_0$, $S_e=\cos W-T_2$, and
$1/(w\sin w)-\cos W=-\cos w+\tfrac{\sqrt\tau}{\sqrt2}\tfrac{\cos^2w}{\sin w}+\cdots$), the $\sqrt\tau$ part is
$-\cos w+T_2$. It cancels because the pole-phase and the bulk saddle carry the \emph{same} leading $\mathrm{lem:cos}$
constant:
\begin{equation}\label{eq:a1}
   \text{\textup{(a1)}}\qquad \cos w_m=\tfrac{\sqrt2}{36}\sqrt\tau\,\sin w+O(\tau^{3/2}),\qquad
   T_2(q_m)=\tfrac{\sqrt2}{36}\sqrt\tau\,\sin w+O(\tau^{3/2}),
\end{equation}
verified $\cos w_m/(\sqrt\tau\sin w),\,T_2/(\sqrt\tau\sin w)\to\tfrac{\sqrt2}{36}=0.03928$ (\texttt{GAPBa\_subatom.py}).
This equality is \emph{structural}: the leading saddle constant is \emph{block-independent} --- both blocks run on
the universal form factor $a_k$ ($g(y)=y/(e^y{-}1)$), and Remark~\ref{rem:blockscope} records the block-dependence
entering only at the \emph{sub}leading $c_1$ (bulk $-\tfrac{17\sqrt2}{36}$, travel $+\tfrac{\sqrt2}{36}$). So
\textup{(a1)} closes the $\sqrt\tau$ cancellation at the Atom~A standard \textup(the leading saddle, for both
blocks\textup), and $P_{12}=O(\tau^{3/2})$ is \emph{proven}. The residual $(\cos w_m-T_2)/\tau^{3/2}\to0.157$ is
bounded. What remains is the single subleading bound
\begin{equation}\label{eq:a2}
   \text{\textup{(a2)}}\qquad |P_{12}|/\tau^{3/2}\ \longrightarrow\ \tfrac1{4\sqrt2}=0.1768\ <\ \tfrac1{\sqrt2}=0.707
   \quad(\text{fourfold margin}),
\end{equation}
the $\tau^{3/2}$-coefficient of the extreme-phase next order. This is the \emph{next-order} van der Corput bound
(Stein \emph{HA} VIII.2 with the leading term subtracted) on the $_0\phi_1$ \eqref{eq:0phi1} --- the exact
q-Bessel analog of Lemma~\ref{lem:extremephase}, the natural $q$-Bessel analogue of the step $V$ carries (now on
the off-diagonal $P_{12}$, not a $V$-block). Its hypotheses are \emph{numerically} verified for the $_0\phi_1$
exactly as Task~F verifies them for $T_2$: simple saddle $k^\ast=w/2$
(\S\ref{rem:qbessel}\,ff.), amplitude $a_k=|d_k|$ of \emph{bounded variation} on the saddle window
($\operatorname{Var}(a)/a(k^\ast)\to2$), and Gaussian-negligible tail ($\sim10^{-6}$); and the value carries a
generous margin, $\sup_m|P_{12}|/\tau^{3/2}=0.187<\tfrac1{\sqrt2}=0.707$ ($3.78\times$, \texttt{GAPBa2\_vdc.py}).
\textbf{This reduces the gate} to (a) $=$ \textup{(a1)} \textup[structural, proven\textup] $+$ \textup{(a2)}
\textup[the subleading residual bound, $3.78\times$ numerical margin\textup]. The \emph{phase} input is closed:
Lemma~\ref{lem:wkbphase} \emph{derives} the constant $\tfrac{\sqrt2}{36}$ as the elementary discrete-WKB
\emph{bulk} phase defect $\tfrac1{24}\sum_n k_n^3$ of the cocycle recursion \eqref{eq:qdiff}, with no
oscillatory-sum saddle and no Airy connection. The \emph{amplitude} input (a2) --- the uniform bound
$|R|=|P_{12}-E|=O(\tau^{5/2})$ --- is the lone outstanding item. Its leading coefficient is the rational
$1891/1296$ (Remark~\ref{rem:qbessel}), so the residual carries no $\sqrt\tau$ oscillatory constant at $O(\tau)$;
but the \emph{uniform} bound is a $q$-Bessel $\nu=\tfrac32$ \emph{turning-point} connection coefficient (the
steepest-descent saddle $k^\ast=w/2$ migrates into the turning region as the phase varies, where the
simple-saddle and bounded-variation hypotheses of the van der Corput estimate above are not established
uniformly), and six independent routes place it at the standing of Lemma~\ref{lem:cos}. \emph{Hence Atom~B reduces
$U$ to a single sharp bound:} $U$ is transcendental \emph{conditional} on $|R|\le K\tau^{5/2}$ --- numerically
certain with a fourfold margin, but unproven --- whereas $V$'s transcendence is \emph{unconditional} (it closes
through the elementary absolute bound Lemma~\ref{lem:T2abs} on a $V$-block; the off-diagonal $P_{12}$ is not a
$V$-block, so $U$'s residual is a \emph{parallel same-class} estimate, not a corollary of $V$).

The cleanest
reading of the gate is through the \emph{ratio} $t_1=P_{12}/S_e$, where the $\mathrm{lem:cos}$ saddle constant
$\tfrac{\sqrt2}{36}$ carried \emph{separately} by $P_{12}$ and by $S_e$ \emph{cancels}, leaving the elementary
leading
\[
   t_1\;\sim\;\tfrac12(w-W)^2\;\longrightarrow\;\tfrac\tau4,\qquad
   s:=\tfrac{q}{1-q}\,t_1\;\longrightarrow\;\tfrac14\;<\;1,
\]
a \emph{fourfold} margin to the gate $s<1$. Numerically $t_1/\tau\to0.25004$ and $s\to0.25004$ monotonically,
with $E/P_{12}\to1$ (the elementary $E$ captures \emph{all} of $P_{12}$ to leading order), and the residual
$R=P_{12}-E$ obeys $R=O(\tau^{5/2})$ with its contribution to $s$ \emph{vanishing} ($\frac{q}{1-q}\frac{R}{S_e}/
\tfrac14\to0$: $0.13,0.05,0.024,\dots\to0$; \texttt{GAPU\_ratio\_check.py}, $80$ poles, \texttt{GAPU\_defect\_trend.py}).
So the gate is carried by the \emph{elementary} leading $t_1\sim\tau/4$ plus the \emph{decay} of $R=P_{12}-E=
O(\tau^{5/2})$ --- the lone outstanding bound \eqref{eq:a2}, numerically certain with a fourfold margin but, per
Remark~\ref{rem:qbessel}, the $q$-Bessel turning-point connection coefficient rather than a closed van der Corput
step. So this
ratio route and the sub-atomisation \textup{(a1)}$+$\textup{(a2)} land at the \emph{same} place: $s\to\tfrac14<1$,
with the lone outstanding input the residual bound \textup{(a2)} of Remark~\ref{rem:qbessel}. \textup(The earlier ``$P_{12}$ directly, bounded coefficient
$A_P<0.706$'' framing is a harder, \emph{unnecessary} route --- the ratio shows the saddle cancels and only the
next-order decay is needed.\textup)

\emph{A tempting shortcut, and why it fails \textup(recorded honestly\textup).} Once the saddle has cancelled,
the ratio $t_1=P_{12}/S_e$ has at the travel poles a $\tau$-expansion with \emph{rational} coefficients,
\begin{equation}\label{eq:t1series}
   \frac{t_1(q_m)}{\tau_m}=\tfrac14+\tfrac3{16}\tau_m+\tfrac{13}{96}\tau_m^2+\tfrac{13}{256}\tau_m^3
   -\tfrac{629}{7680}\tau_m^4-\cdots,
\end{equation}
verified to $16$ orders (integer-relation detection on $1004$-digit data). To low order the coefficients are small
($|c_k|\sim0.1$--$0.2$), and summing the first ten at the largest pole $\tau_1=0.0905$ reproduces $t_1(q_1)$ to
$\sim10^{-11}$, suggesting a \emph{convergent} series whose summation would give $s<1$ with no error term.
\emph{This is an artifact of stopping early.} Extending to $25$ orders gives $|c_k|^{1/k}\sim0.068\,k+0.38$
(slope $>0$): the series is in fact \emph{divergent} (Gevrey-$1$, zero radius of convergence), matching the
divergence of the amplitude data $a_n$ and the off-axis Borel branch cut of the same $q$-series. So $t_1/\tau$
\emph{cannot} be summed, and the residual is \emph{not} the mild ``radius $>\tau_1$'' statement --- it remains the
same asymptotic tail, the turning-point connection coefficient of Remark~\ref{rem:qbessel}. The convergent-series
shortcut is therefore \emph{withdrawn}; the honest residual is the single uniform bound $|R|\le K\tau^{5/2}$.
\end{remark}
